\documentclass{article}
\usepackage{anyfontsize}
\usepackage[UTF8]{ctex} % 如果需要支持中文
\usepackage{fontspec} 

\usepackage[a4paper, left=0.1in, right=0.1in, top=0.05in, bottom=0.05in]{geometry} % 设置四边页边距为0.1英寸{geometry} % 设置页边距为0.5英寸
\usepackage{enumitem} % 用于自定义列表
\usepackage{amsmath}  % 数学公式支持

\setlength{\parindent}{0pt} % 不缩进
\usepackage{etoolbox} % 用于列表操作
%调整类代码
\usepackage{comment}
\usepackage{tocloft} % 用于定制目录样式
\usepackage{hyperref} %不需要目录盒子注释这
\usepackage{ifthen}
\usepackage{tikz}
\usepackage{titletoc}
\usepackage{graphicx}
\usepackage{amssymb}  % 提供数学符号，包括\therefore
%目录设定
%快捷键 CMD + / 注释
%  \renewcommand{\cftdot}{} % 隐藏目录中的页码
%  \cftpagenumbersoff{section}
%  \cftpagenumbersoff{subsection}
%  \makeatletter
%  \renewcommand{\@pnumwidth}{0em}  % 设置页码宽度为0
%  \renewcommand{\@tocrmarg}{0em}   % 设置右边距为0
%  \makeatother
 
%\setCJKmainfont[
 %   Path = C:/USERS/11252/APPDATA/LOCAL/MICROSOFT/WINDOWS/FONTS/,
 %   UprightFont = {SourceHanSerifCN-Light-5.otf},
 %   BoldFont = {SourceHanSerifCN-Medium-6.otf},
%    Scale=1
%]{SourceHanSerifCN}

%\setmainfont{SourceHanSerif}  % 设置英文字体

% 处理冲突的命令
\let\oldbackepsilon\backepsilon
\let\backepsilon\relax
\let\oldeth\eth
\let\eth\relax
\let\olddigamma\digamma
\let\digamma\relax

\usepackage{float}
\usepackage[a4paper, left=0.1in, right=0.1in, top=0.05in, bottom=0.05in]{geometry} % 设置四边页边距为0.1英寸{geometry} % 设置页边距为0.5英寸
\usepackage{enumitem} % 用于自定义列表
\usepackage{amsmath}  % 数学公式支持

\setlength{\parindent}{0pt} % 不缩进
\usepackage{etoolbox} % 用于列表操作
%调整类代码
\usepackage{comment}
\usepackage{tocloft} % 用于定制目录 样式
\usepackage{hyperref} %不需要目录盒子注释这
\usepackage{ifthen}
\usepackage{tikz}
\usepackage{titletoc}
\usepackage{graphicx}
\usepackage{amssymb}  % 提供数学符号，包括\therefore
%目录设定
%快捷键 CMD + / 注释
%  \renewcommand{\cftdot}{} % 隐藏目录中的页码
%  \cftpagenumbersoff{section}
%  \cftpagenumbersoff{subsection}
%  \makeatletter
%  \renewcommand{\@pnumwidth}{0em}  % 设置页码宽度为0
%  \renewcommand{\@tocrmarg}{0em}   % 设置右边距为0
%  \makeatother
 


\usepackage{environ}

\newif\ifshowcontent  % 定义一个条件判断变量
\showcontenttrue    % 设置为 false，隐藏所有 question 环境的内容

\newcounter{question} % 题目计数器
\setcounter{question}{0}
\NewEnviron{question}{%
    \ifshowcontent
        \par\stepcounter{question}\textbf{\thequestion.} \BODY\par % 如果显示内容，则输出
    \fi
}

\NewEnviron{choices}{%
    \ifshowcontent
        \begin{enumerate}[label=\Alph*., leftmargin=*]
            \BODY
        \end{enumerate}\par
    \fi
}

\NewEnviron{correctanswer}{%
    \ifshowcontent
        \textbf{答案:}\BODY\par
    \fi
}



\newenvironment{explanation}
{\textbf{解析:}\par}
{\par}



% 新增考察方面命令
\newcommand{\checklist}[1]{
    \textbf{考察:} #1 \par % 在题目下方显示考察方面
    \addcontentsline{toc}{subsection}{题号\thequestion 考察: #1} % 将考察内容添加到目录
    \vspace{2em} % 添加1em的垂直空间
}

\graphicspath{{images/}} %统一


\begin{document}
 %\fontsize{22}{31}\selectfont
 \tableofcontents % 添加目录
\newpage
%\pagestyle{empty}




\section{考点1:信用风险建模评估}
\setcounter{question}{0}

\begin{question}
    Which of the following statements regarding a banking credit analyst's skills is most likely correct?
\end{question}

\begin{choices}
    \item High earnings quality suggests that the bank is profitable
    \item Peer analysis is facilitated by the standardized nature of financial performance measures
    \item Although qualitative analytical skills are required, quantitative analytical skills are more important
    \item In analyzing an unfamiliar banking sector, an analyst should start by performing detailed reviews of the major banks
\end{choices}

\begin{correctanswer}
    B
\end{correctanswer}

\begin{explanation}
    \textbf{A选项错误。}
    \begin{itemize}
        \item 高盈利质量不一定直接等同于银行盈利能力。
        \item 盈利质量衡量的是利润的可持续性和可靠性（如是否依赖一次性收益或会计调整），而盈利（profitability）指整体利润规模。
    \end{itemize}
    
    \textbf{B选项正确。}
    \begin{itemize}
        \item 同业比较需要统一会计标准（如IFRS或GAAP）和财务比率（如ROA、NIM等），才能有效横向对比不同银行的信用风险。
        \item 所以标准化性质使得同业分析更便利。
    \end{itemize}
    
    \textbf{C选项错误。}
    \begin{itemize}
        \item 定量和定性分析技能同样重要，它们服务于不同（但相关）的目的。
        \item 定性技能有助于确定实体偿还债务的意愿，而定量技能有助于确定实体偿还债务的能力。
    \end{itemize}
    
    \textbf{D选项错误。}
    \begin{itemize}
        \item 分析陌生银行业时应先宏观（行业环境、监管框架）→再微观（个体银行财务指标）。
        \item 直接详细分析大银行可能忽视系统性风险。
    \end{itemize}
\end{explanation}

\checklist{商业银行信用分析师的技能要求}



\begin{question}
    An exposure has a stressed default probability of 4.0\% and loss given default of 50.0\%. The portfolio with a value of \$1 million. The standard deviation of the LGD is 25.0\%. What is the stressed expected loss?
\end{question}

\begin{choices}
    \item 0.02 million
    \item 0.0125 million
    \item 0.04 million
    \item 0.05 million
\end{choices}

\begin{correctanswer}
    A
\end{correctanswer}

\begin{explanation}
\textbf{预期损失（Stressed Expected Loss, EL）计算公式} 
    \[
    \text{EL} = \text{EAD} \times \text{PD} \times \text{LGD}
    \]
    
    \textbf{题设数据：}
    \begin{itemize}
      \item 风险敞口（EAD）：\$1{,}000{,}000
      \item 违约概率（PD）：4\% = 0.04
      \item 违约损失率（LGD）：50\% = 0.5
    \end{itemize}
    
    \textbf{计算：}
    \[
    \text{Stressed EL} = 1{,}000{,}000 \times 0.04 \times 0.5 = 20{,}000 = 0.02 \text{ million}
    \]
    
    \textbf{干扰项说明：}题目中提到的 LGD 标准差（25\%）。用于 非预期损失（UL） 或信用 VaR 的估计，与本题无关。
    
    \end{explanation}

\checklist{预期损失计算，注意排除题干干扰数据}


\begin{question}
    Which of the following statements regarding the role of a corporate credit analyst is most likely correct?
\end{question}

\begin{choices}
    \item Earnings analysis is by far the most important analyst task
    \item The larger the size of the firm, the lower the cost of analysis
    \item Analysts are generally required to cover multiple industry areas given the huge diversity among corporations
    \item The smaller the firm, the lower the cost of analysis
\end{choices}

\begin{correctanswer}
    B
\end{correctanswer}

\begin{explanation}
    \textbf{A选项错误。}
    \begin{itemize}
        \item 信用分析是综合性工作，需平衡财务分析（如现金流、杠杆比率）与非财务因素（如管理层能力、行业风险）。
        \item 盈利分析仅是其中一环，而非“最重要”任务。
    \end{itemize}
    
    \textbf{B选项正确。}
    \begin{itemize}
        \item 大型企业通常有更透明的财务报表和更规范的披露流程，减少了信息收集成本。
        \item 此外，大型企业的历史数据更易获取，便于标准化分析。
    \end{itemize}
    
    \textbf{C选项错误。}
    \begin{itemize}
        \item 分析师通常专注于特定行业（如 banking、能源），以积累深度专业知识，而非分散精力跨行业。
    \end{itemize}
    
    \textbf{D选项错误。}
    \begin{itemize}
        \item 中小企业因信息不透明、数据缺失，分析成本反而更高。需依赖更多实地调研和非标准化数据。
    \end{itemize}
\end{explanation}

\checklist{企业信用分析的工作特征与现实挑战（注意中小企业分析成本更高）}


\begin{question}
    Which of the following is true about predatory lending and predatory borrowing?
\end{question}

\begin{choices}
    \item Both underprovide credit
    \item Both overprovide credit
    \item Predatory lending underprovides credit and predatory borrowing overprovides credit
    \item Predatory lending overprovides credit and predatory borrowing underprovides credit
\end{choices}

\begin{correctanswer}
    B
\end{correctanswer}

\begin{explanation}
    \textbf{A选项错误。}
    \begin{itemize}
        \item 掠夺性放贷和掠夺性借贷都是借款方过度授信。
    \end{itemize}
    
    \textbf{B选项正确。}
    \begin{itemize}
        \item 掠夺性放贷（Predatory Lending）贷款方主动诱导或强制借款人接受不合理条款（如超高利率、隐藏费用），利用信息不对称谋利。
        \item 掠夺性借贷（Predatory Borrowing）是借款方通过欺诈或虚假信息获取贷款，无真实还款意愿。
        \item 掠夺性借贷的欺诈性本质会引发后续信贷紧缩，导致金融机构收紧风控标准，减少对市场的信用供给。
    \end{itemize}
    
    \textbf{C选项错误。}
    \begin{itemize}
        \item 掠夺性借贷误导了贷款人对借款人信用状况的判断，导致授予了过量信贷。
    \end{itemize}
    
    \textbf{D选项错误。}
    \begin{itemize}
        \item 掠夺性借贷会使银行授信过度，会引发市场信用收缩，最终成为授信不足的诱因。
        \item 注意区分题目所问。
    \end{itemize}
\end{explanation}

\checklist{考察掠夺性借贷和掠夺性放贷（注意区分问行为还是结果）}


\begin{question}
    Which of the following model calculates the change in portfolio value due to rating migration of the underlying instruments?
\end{question}

\begin{choices}
    \item Credit Risk+
    \item Credit Metrics
    \item KMV
    \item Both a and c above are true
\end{choices}

\begin{correctanswer}
    B
\end{correctanswer}

\begin{explanation}
    \textbf{A选项错误。}
    \begin{itemize}
        \item Credit Risk+模型以违约率为基础的统计学模型，将违约视作泊松过程。
        \item 模拟的是信贷损失分布，重点是违约频率，不考虑评级变化。
    \end{itemize}
    
    \textbf{B选项正确。}
    \begin{itemize}
        \item CreditMetrics模型基于信用评级迁移矩阵，计算因评级变动导致的组合价值重估。
    \end{itemize}
    
    \textbf{C选项错误。}
    \begin{itemize}
        \item KMV模型基于Merton结构模型，预测违约概率，不直接处理评级迁移。
    \end{itemize}
    
    \textbf{D选项错误。}
    \begin{itemize}
        \item Credit Risk+和KMV均不处理评级迁移，故排除。
    \end{itemize}
\end{explanation}

\checklist{理解不同信用风险计量模型的侧重点}

\begin{question}
    Lisa Chen is a new loan officer at Pacific Bank. Her manager requires her to make five commercial loans this month. Chen hears about a local entrepreneur with a good reputation. She meets with the entrepreneur, David Wang, and is impressed by his business acumen. They discuss a loan for Wang's new project, and Chen agrees it's a promising opportunity. She convinces the loan committee that Pacific Bank is fortunate to work with someone of Wang's reputation. This is an example of:
\end{question}

\begin{choices}
    \item Historical analysis technique
    \item Qualitative analysis technique
    \item Quantitative analysis technique
    \item Extrapolation analysis technique
\end{choices}

\begin{correctanswer}
    B
\end{correctanswer}

\begin{explanation}
    \textbf{A选项错误。}
    \begin{itemize}
        \item 历史分析技术常指借助过去的财务数据、违约记录、历史现金流等构建的分析体系。
        \item 案例未涉及借贷人过往的信用记录，故排除。
    \end{itemize}
    
    \textbf{B选项正确。}
    \begin{itemize}
        \item 名义贷款（Name lending）是一种定性分析技术，有时用来代替财务分析。
        \item 定性分析通过非数值因素（如管理层能力、行业前景、声誉）评估信用风险。
        \item 案例中Chen依赖Wang的声誉和主观判断，符合定性分析。
    \end{itemize}
    
    \textbf{C选项错误。}
    \begin{itemize}
        \item 定量分析技术指使用数学模型和数值数据（如财务比率、信用评分）评估风险。
        \item 案例中未涉及量化指标。
    \end{itemize}
    
    \textbf{D选项错误。}
    \begin{itemize}
        \item 外推分析技术通过历史趋势外推未来风险（如时间序列预测）。
        \item 案例中未提及趋势外推。
    \end{itemize}
\end{explanation}

\checklist{信用风险分析方法的分类（定量vs定性）}


\begin{question}
    A bank has an outstanding trade with one of its counterparties with an exposure of \$500,000 and a recovery rate of 70\%. The bank estimated that there is a 2\% probability that the counterparty will default on its obligations. What is the bank's expected loss?
\end{question}

\begin{choices}
    \item \$3,000
    \item \$7,000
    \item \$10,000
    \item \$150,000
\end{choices}

\begin{correctanswer}
    A
\end{correctanswer}

\begin{samepage}
\begin{explanation}
  \textbf{预期损失（Stressed Expected Loss, EL）计算公式：} 
    \[
    \text{EL} = \text{EAD} \times \text{PD} \times \text{LGD}
    \]

    \textbf{题设数据：}
    \begin{itemize}
      \item 风险敞口（EAD）：\$500{,}000
      \item 违约概率（PD）：2\% 
    \end{itemize}

     
       \textbf{步骤 1：计算违约损失率（LGD）}
        \begin{itemize}
            \item 回收率 = 70\%
            \item 违约损失率（LGD） =1 - 70\% = 30\%
        \end{itemize}
        
    
      \textbf{步骤 2：计算预期损失}
        \begin{itemize}
            \item 预期损失 =\$500{,}000 × 2\% × 30\% = \$3,000
        \end{itemize}
    
    \end{explanation}
    
\checklist{预期损失的计算,注意辨别回收率和违约损失率}
\end{samepage}




\begin{question}
    The KMV model measures the normalized "distance to default". How is this defined?
\end{question}

\begin{choices}
    \item (Expected Assets - Default Point) / (Volatility of assets)
    \item Equity / (Volatility of equity)
    \item Probability of stock price falling below a threshold
    \item Leverage x Stock Price Volatility
\end{choices}

\begin{correctanswer}
    A
\end{correctanswer}

\begin{explanation}
    \textbf{A选项正确。}
    \begin{itemize}
        \item 违约距离（Distance to Default）的定义是：（预期资产 - 加权债务）/（资产波动率）。
        \item 这是一个标准化的违约衡量指标，因此可以用于比较不同公司。
    
    \end{itemize}
    
    \textbf{B选项错误。}
    \begin{itemize}
        \item KMV模型基于资产价值而非股权价值。
        \item 股权波动率因杠杆效益通常高于资产波动率，不能直接替代。
    \end{itemize}
    
    \textbf{C选项错误。}
    \begin{itemize}
        \item 描述的是股价触发的违约概率，属于后续转化步骤。
        \item 违约距离是一个标准化的距离指标，而不是概率。
    \end{itemize}
    
    \textbf{D选项错误。}
    \begin{itemize}
        \item 杠杆与股价波动率的乘积是市场风险指标，与KMV的资产价值驱动逻辑无关。
    \end{itemize}
\end{explanation}

\checklist{KMV模型中的违约距离是（预期资产-加权债务）/资产波动率}

\begin{question}
    An analyst is using Moody's KMV model to estimate the distance to default of a large public firm, Sports Inc., a firm that designs, manufactures and sells athletic shoes. The firm's capital structure consists of USD 40 million in short-term debt, USD 20 million in long-term debt, and there are one million shares of stock currently trading at USD 10 per share. The asset volatility is 20\% per year. What is the normalized distance to default for Sports Inc.?
\end{question}

\begin{choices}
    \item 0.714
    \item 1.430
    \item 2.240
    \item 5.000
\end{choices}

\begin{correctanswer}
    B
\end{correctanswer}

\begin{explanation}
    \textbf{根据KMV模型，标准化违约距离（Distance to Default）计算公式}
    \begin{equation}
        \text{Distance to Default (DD)} = \displaystyle \frac{V_A - D}{\sigma_A \cdot V_A}
    \end{equation}

\textbf{步骤1：计算公司资产价值$V_A$}
   
    \begin{itemize}
        \item 总负债 = 短期负债 + 长期负债  
        \item 股权价值$V_A$＝股权市值＋总负债
        \item 则$V_A$＝1 million shares × \$10/share + (40 million + 20 million ) = 70 million
     \end{itemize}

\textbf{步骤2：计算违约点$D$（Default Point）}

    \begin{itemize}
         \item $D = \text{短期负债} + \text{长期负债} \times 0.5 = 40\,\text{million} + 20\,\text{million} \times 0.5 = 50\,\text{million}$
      
    \end{itemize}

\textbf{步骤3：计算违约距离（DD）}
  

    \begin{itemize}
        \item 违约距离的公式为：$\displaystyle \text{DD} = \frac{V_A - D}{\sigma_A \cdot V_A}$
        \item 将数据代入：$\displaystyle \text{DD} = \frac{70\,\text{million} - 50\,\text{million}}{0.2 \times 70\,\text{million}}$
        \item 得：$ \text{DD} = \frac{20\,\text{million}}{14\,\text{million}} \approx 1.430$
    \end{itemize}  
  
\end{explanation}

\checklist{熟悉KMV模型中的违约距离计算公式（违约点中长期负债只需算一半）}

\begin{question}
    David Wang is a senior analyst at an investment bank preparing research reports on the mining sector. Wang noted that one of the companies he follows recently increased risk to the firm's assets, which he expects will benefit equity holders to the detriment of debt holders. Which of the following concepts best describes the scenario in Wang's analysis?
\end{question}

\begin{choices}
    \item Coordination failures
    \item Adverse selection
    \item Risk shifting
    \item Principal-agent problem
\end{choices}

\begin{correctanswer}
    C
\end{correctanswer}

\begin{explanation}
    \textbf{A选项错误。}
    \begin{itemize}
        \item 协调失败（Coordination failures）指多方因缺乏协调导致集体非最优结果（如银行挤兑）。
        \item 案例未涉及多方协作问题。
    \end{itemize}
    
    \textbf{B选项错误。}
    \begin{itemize}
        \item 逆向选择（Adverse selection）指信息不对称下高风险主体更倾向于交易（如次贷危机）。
        \item 案例中无信息不对称或交易选择行为。
    \end{itemize}
    
    \textbf{C选项正确。}
    \begin{itemize}
        \item 风险转移（Risk shifting）指公司股东或管理层将更多风险“转嫁”给债权人，从而博取更高的潜在回报。
        \item 在本例中，公司资产风险的增加会使股权持有人受益，而损害债务持有人的利益，因为债务持有人有固定收益但承担了更大的损失风险。
    \end{itemize}
    
    \textbf{D选项错误。}
    \begin{itemize}
        \item 委托-代理问题（Principal-agent problem）常发生在管理层（代理人）与股东（委托人）利益冲突时
        \item 案例中矛盾主体为股东与债权人，而非管理层。
    \end{itemize}
\end{explanation}

\checklist{信用风险中风险转移（Risk Shifting）概念的掌握（较简单的情景题）}



\newpage



\section{考点2:信用评分与评级}
\setcounter{question}{0}

\begin{question}
    Michael Chen, FRM, is a rating agency analyst who is currently performing financial statement analysis on a major bank. Which of the following financial statements would be least useful for bank credit analysis?
\end{question}

\begin{choices}
    \item Balance sheet
    \item Income statement
    \item Statement of cash flows
    \item Statement of changes in capital funds
\end{choices}

\begin{correctanswer}
    C
\end{correctanswer}

\begin{explanation}
    \textbf{A选项错误。}
    \begin{itemize}
        \item Balance sheet（资产负债表）映银行的资产质量（如贷款组合）、资本充足性（资本比率）和流动性（现金及等价物）
        \item 信用分析的核心依据。
    \end{itemize}
    
    \textbf{B选项错误。}
    \begin{itemize}
        \item  Income statement（利润表）展示银行的盈利能力（如净利息收入、不良贷款拨备）
        \item 影响银行偿债能力评估。
    \end{itemize}
    
    \textbf{C选项正确。}
    \begin{itemize}
        \item  Statement of cash flows（现金流量表）展示企业在某一期间的经营、投资、筹资活动中现金流入和流出。
        \item  对制造业或传统企业极为重要，但在银行信用分析中价值较低：
        \begin{enumerate}
            \item 银行本身的“经营活动”就是存放款与金融交易，现金流量反映不出真实业务结构
            \item 银行业的资产高度流动化，且现金变动往往不等于“流动性风险”
            \item 银行获取资金依赖于市场融资能力和资本充足，而非经营现金净流入
        \end{enumerate}
    \end{itemize}
    
    \textbf{D选项错误。}
    \begin{itemize}
        \item  Statement of changes in capital funds（资本变动表）显示了银行资本的变化情况
        \item 影响银行资本充足性和恢复能力评估
    \end{itemize}
\end{explanation}

\checklist{银行信用分析中各报表的重要性判断（区别现金流量表对各行业信用分析的重要性）}

\begin{question}
    Each of the following items represents an example of qualitative information that would ideally be captured in assessing default probability except:
\end{question}

\begin{choices}
    \item Management's education and experience
    \item Internal controls associated with financial reporting
    \item Diversification of products and customers locally and globally
    \item Trends in throughput and other operational efficiency metrics
\end{choices}

\begin{correctanswer}
    D
\end{correctanswer}

\begin{explanation}
    \textbf{A选项错误。}
    \begin{itemize}
        \item 管理层能力是典型的主观评估因素，直接影响公司战略和风险应对能力
    
    \end{itemize}
    
    \textbf{B选项错误。}
    \begin{itemize}
        \item 财务报告内部控制有效性需通过审计意见或流程描述评估，属于非量化指标
    \end{itemize}
    
    \textbf{C选项错误。}
    \begin{itemize}
        \item 业务多元化程度通常通过描述性分析（如市场份额分布）判断
    \end{itemize}
    
    \textbf{D选项正确。}
    \begin{itemize}
        \item 吞吐量（如单位时间产量）、生产效率等是可量化的运营数据
        \item 直接对应数值分析，属于定量信息
    \end{itemize}
\end{explanation}

\checklist{定性与定量信息的甄别（注意题目需求是定性信息）}


\begin{question}
    Golin and Delhaise divide credit analysis into four areas according to borrower type:
    \begin{enumerate}
        \item Consumer credit analysis: evaluation of individual consumers' creditworthiness
        \item Corporate credit analysis: evaluation of nonfinancial companies
        \item Financial institution credit analysis: evaluation of financial companies
        \item Sovereign/municipal credit analysis: evaluation of credit risk of nations and governments
    \end{enumerate}
    According to Golin and Delhaise, each of the following is true about key features of credit analysis with respect to borrower type, except which is not true?
\end{question}

\begin{choices}
    \item Individuals (consumers): Credit analysis is amenable to automation and the use of scoring models and statistical tools to correlate risk to limited number of variables
    \item Non-financial corporations: Compared to consumers, tends to be more detailed and "hands-on" (i.e. less automated); key variables are likely to include liquidity, cash flow, near-term earnings capacity and profitability, solvency or capital position
    \item Financial Companies: In contrast to corporate (non-financial) credit analysis, qualitative analysis and asset quality are not important, but cash flow is a highly important (a "key indicator")
    \item Sovereigns: Includes analysis of country risk, which is primarily political dynamics and state of the economy; and systematic risk, which includes the regulatory regime and the financial system
\end{choices}

\begin{correctanswer}
    C
\end{correctanswer}

\begin{explanation}
    \textbf{A选项正确，但不选。}
    \begin{itemize}
        \item 消费者信用分析的自动化程度高
        \item 依赖信用评分模型（如FICO）和少量核心变量（收入、负债比等）
    \end{itemize}
    
    \textbf{B选项正确，但不选。}
    \begin{itemize}
        \item 非金融企业分析更复杂且依赖人工判断
        \item 需结合流动性、现金流、盈利能力和资本充足性等定量指标
    \end{itemize}
    
    \textbf{C选项错误。}
    \begin{itemize}
        \item 金融机构信用分析中定性分析（管理层、治理、风险文化）和资产质量至关重要
        \item 现金流并非唯一关键指标（如现金流结构复杂的金融机构）
    \end{itemize}
    
    \textbf{D选项正确，但不选。}
    \begin{itemize}
        \item 主权风险分析涵盖国家风险分析（政治经济稳定性）和系统性风险（如金融监管体系）
    \end{itemize}
\end{explanation}

\checklist{信用分析分类框架（区分不同借款人信用分析的侧重点）}


\newpage


\section{考点3:违约可能性评估}
\setcounter{question}{0}
\begin{question}
    Which of the following statements is (are) not true regarding a credit default swap (CDS)?
    \begin{enumerate}
        \item The purchaser of a CDS seeks credit protection
        \item Swap partners make payments based on a benchmark credit spread on the reference obligation
        \item Physical settlement provides the CDS buyer the par value of the reference obligation
        \item A CDS has the same payoffs as a credit default put with installment payments
    \end{enumerate}
\end{question}

\begin{choices}
    \item II only
    \item I and III
    \item III only
    \item II and IV
\end{choices}

\begin{correctanswer}
    A
\end{correctanswer}

\begin{explanation}
    \textbf{A选项正确。}
    \begin{itemize}
        \item 陈述II不正确，因为：
        \begin{enumerate}
            \item CDS买方在互换期限内或发生信用事件之前，向CDS卖方支付固定费用
            \item 支付不是基于参考债务的基准信用利差
            \item 支付是预先确定的固定金额
        \end{enumerate}
    \end{itemize}
    
    \textbf{B选项错误。}
    \begin{itemize}
        \item 陈述I正确：CDS买方确实寻求信用保护
        \item 陈述III正确：实物结算确实给CDS买方提供参考债务的面值
    \end{itemize}
    
    \textbf{C选项错误。}
    \begin{itemize}
        \item 陈述III正确：实物结算确实给CDS买方提供参考债务的面值
    \end{itemize}
    
    \textbf{D选项错误。}
    \begin{itemize}
        \item 陈述II不正确：支付不是基于基准信用利差
        \item 陈述IV正确：CDS确实与分期付款的信用违约看跌期权有相同的收益
    \end{itemize}
\end{explanation}

\checklist{CDS买方支付固定费用，而不是基于基准信用利差}

\begin{question}
    Which of the following statements is most likely correct about recovery rates?
\end{question}

\begin{choices}
    \item Recovery rates vary inversely with capital structure seniority
    \item Historical recovery rates are fairly constant across industries
    \item Recovery rates are highest during economic downturns
    \item Actual recovery rates may differ substantially from settled recovery rates
\end{choices}

\begin{correctanswer}
    D
\end{correctanswer}

\begin{explanation}
    \textbf{A选项错误。}
    \begin{itemize}
        \item 回收率与资本结构优先级成正比，而不是反比：
        \begin{enumerate}
            \item 优先级越高的债务，回收率越高
            \item 优先级越低的债务，回收率越低
        \end{enumerate}
    \end{itemize}
    
    \textbf{B选项错误。}
    \begin{itemize}
        \item 历史回收率在不同行业之间存在显著差异：
        \begin{enumerate}
            \item 不同行业的资产特征不同
            \item 不同行业的破产程序不同
            \item 不同行业的经济周期敏感性不同
        \end{enumerate}
    \end{itemize}
    
    \textbf{C选项错误。}
    \begin{itemize}
        \item 回收率在经济衰退期间最低，而不是最高：
        \begin{enumerate}
            \item 经济衰退时资产价值下降
            \item 经济衰退时破产程序可能更复杂
            \item 经济衰退时市场流动性可能更差
        \end{enumerate}
    \end{itemize}
    
    \textbf{D选项正确。}
    \begin{itemize}
        \item 实际回收率可能与结算回收率有显著差异，因为：
        \begin{enumerate}
            \item 结算回收发生在信用事件后不久（如CDS拍卖或违约债券出售）
            \item 实际回收可能发生在多年后，基于破产解决结果
            \item 破产程序可能持续很长时间
            \item 最终回收金额可能受到多种因素影响
        \end{enumerate}
    \end{itemize}
\end{explanation}

\checklist{实际回收率可能与结算回收率有显著差异，因为结算回收发生在信用事件后不久，而实际回收可能发生在多年后}


\begin{question}
    A 2-year credit default swap (CDS) specifying physical delivery defaults at the end of two years. If the reference asset is a \$200 million, 8.0\% ABC corporate bond, and the CDS spread is 125 basis points, the buyer of the CDS will:
\end{question}

\begin{choices}
    \item Receive payments of 800 basis points for the next two years
    \item Receive a payment of \$167.5 million
    \item Deliver the bond and receive a payment of \$200 million
    \item Continue to receive payments of 675 basis points for the next two years
\end{choices}

\begin{correctanswer}
    C
\end{correctanswer}

\begin{explanation}
    \textbf{A选项错误。}
    \begin{itemize}
        \item CDS买方不是接收支付，而是支付固定费用（CDS spread）给卖方
        \item 800个基点是债券的票面利率，不是CDS的支付
    \end{itemize}
    
    \textbf{B选项错误。}
    \begin{itemize}
        \item \$167.5 million不是正确的支付金额
        \item 在实物交割的CDS中，买方会收到参考债务的面值（\$200 million）
    \end{itemize}
    
    \textbf{C选项正确。}
    \begin{itemize}
        \item 在实物交割的CDS中：
        \begin{enumerate}
            \item CDS买方将参考债务（ABC公司债券）交付给卖方
            \item CDS买方收到参考债务的面值（\$200 million）
            \item 这是实物交割CDS的标准结算方式
        \end{enumerate}
    \end{itemize}
    
    \textbf{D选项错误。}
    \begin{itemize}
        \item CDS买方不是接收支付，而是支付固定费用（CDS spread）给卖方
        \item 675个基点不是正确的CDS spread（实际是125个基点）
    \end{itemize}
\end{explanation}

\checklist{实物交割的CDS中，买方将参考债务交付给卖方，并收到参考债务的面值}


\begin{question}
    Which of the following statements about credit default swaps is most accurate?
\end{question}

\begin{choices}
    \item CDSs transfer credit risk and market risk from the protection buyer to the protection seller
    \item CDSs transfer credit risk from the protection buyer to the issuer of the underlying credit
    \item Physical settlement requires knowledge of the post-default market price
    \item Cash settlement avoids the problem of a delivery squeeze
\end{choices}

\begin{correctanswer}
    D
\end{correctanswer}

\begin{explanation}
    \textbf{A选项错误。}
    \begin{itemize}
        \item CDS只转移信用风险，不转移市场风险：
        \begin{enumerate}
            \item CDS主要针对信用事件提供保护
            \item 市场风险（如利率风险、汇率风险等）不在CDS的覆盖范围内
        \end{enumerate}
    \end{itemize}
    
    \textbf{B选项错误。}
    \begin{itemize}
        \item CDS将信用风险从保护买方转移给保护卖方，而不是基础信用的发行人：
        \begin{enumerate}
            \item 保护卖方承担信用风险
            \item 基础信用的发行人仍然是信用风险的源头
        \end{enumerate}
    \end{itemize}
    
    \textbf{C选项错误。}
    \begin{itemize}
        \item 实物结算不需要知道违约后的市场价格：
        \begin{enumerate}
            \item 实物结算基于参考债务的面值
            \item 市场价格与实物结算无关
        \end{enumerate}
    \end{itemize}
    
    \textbf{D选项正确。}
    \begin{itemize}
        \item 现金结算避免了交割挤压问题，因为：
        \begin{enumerate}
            \item 没有实际的证券交易
            \item 保护买方不需要在公开市场购买参考实体
            \item 不会因为大量购买导致价格上升
        \end{enumerate}
    \end{itemize}
\end{explanation}

\checklist{现金结算避免了交割挤压问题，因为不需要实际的证券交易}


\begin{question}
    Economic capital calculations for credit risk assume a recovery rate (defined as 1-loss rate). Recovery rates are dependent on the business model of the underlying counterparty and its asset volatility in value and size. Under normal anticipated circumstances which of the following types of companies will have the highest recovery rate?
\end{question}

\begin{choices}
    \item An internet merchant of designer clothes
    \item A hedge fund
    \item An asset intensive manufacturing company
    \item A commodities trader
\end{choices}

\begin{correctanswer}
    C
\end{correctanswer}

\begin{explanation}
    \textbf{A选项错误。}
    \begin{itemize}
        \item 网上设计师服装零售商通常：
        \begin{enumerate}
            \item 有形资产较少
            \item 主要资产是库存和应收账款
            \item 在违约时难以变现
        \end{enumerate}
    \end{itemize}
    
    \textbf{B选项错误。}
    \begin{itemize}
        \item 对冲基金通常：
        \begin{enumerate}
            \item 有形资产很少
            \item 主要资产是金融资产
            \item 资产价值波动较大
        \end{enumerate}
    \end{itemize}
    
    \textbf{C选项正确。}
    \begin{itemize}
        \item 资产密集型制造公司通常：
        \begin{enumerate}
            \item 拥有大量有形资产（如厂房、设备）
            \item 资产价值相对稳定
            \item 在违约时容易估值和变现
        \end{enumerate}
        \item 例如，公用事业公司有很高的回收率，因为它们拥有大量有形资产，如发电厂
    \end{itemize}
    
    \textbf{D选项错误。}
    \begin{itemize}
        \item 商品交易商通常：
        \begin{enumerate}
            \item 有形资产较少
            \item 主要资产是商品库存
            \item 资产价值波动较大
        \end{enumerate}
    \end{itemize}
\end{explanation}

\checklist{资产密集型制造公司通常拥有大量有形资产，在违约时容易估值和变现，因此回收率最高}


\begin{question}
    An analyst has gathered the following information about ABC Inc. and DEF Inc. The respective credit ratings are AA and BBB with 1-year CDS spreads of 200 and 300 basis points each. The associated probabilities of default based on published reports are 10\% and 20\%, respectively. Which of the following statements about the recovery rates is most likely correct?
\end{question}

\begin{choices}
    \item The market implied recovery rates are equal
    \item The market implied recovery rate is higher for ABC
    \item The market implied recovery rate is lower for ABC
    \item The loss given default is higher for DEF
\end{choices}

\begin{correctanswer}
    C
\end{correctanswer}

\begin{explanation}
    \textbf{A选项错误。}
    \begin{itemize}
        \item 市场隐含回收率不相等：
        \begin{enumerate}
            \item ABC的回收率 = 80\%
            \item DEF的回收率 = 85\%
        \end{enumerate}
    \end{itemize}
    
    \textbf{B选项错误。}
    \begin{itemize}
        \item ABC的市场隐含回收率（80\%）低于DEF（85\%）
        \item 计算过程：
        \begin{enumerate}
            \item 信用利差 ≈ (1 - 回收率) × 违约概率
            \item ABC: 200 bps = (1 - RR) × 10\%, 所以RR = 80\%
            \item DEF: 300 bps = (1 - RR) × 20\%, 所以RR = 85\%
        \end{enumerate}
    \end{itemize}
    
    \textbf{C选项正确。}
    \begin{itemize}
        \item ABC的市场隐含回收率（80\%）确实低于DEF（85\%）
        \item 使用违约损失率（LGD）表示：
        \begin{enumerate}
            \item ABC的LGD = 20\%
            \item DEF的LGD = 15\%
        \end{enumerate}
    \end{itemize}
    
    \textbf{D选项错误。}
    \begin{itemize}
        \item DEF的违约损失率（15\%）低于ABC（20\%）
        \item 违约损失率 = 1 - 回收率
    \end{itemize}
\end{explanation}

\checklist{市场隐含回收率可以通过信用利差和违约概率计算得出，ABC的回收率（80\%）低于DEF（85\%）}

\begin{question}
    The buyer of a credit-default swap (CDS):
\end{question}

\begin{choices}
    \item Has no risk exposure to the combined holdings of the CDS and reference obligation
    \item Is subject to the credit risk of the seller
    \item Must place collateral with the seller of the CDS to ensure performance
    \item Can surrender the CDS for the value of the accumulated payments
\end{choices}

\begin{correctanswer}
    B
\end{correctanswer}

\begin{explanation}
    \textbf{A选项错误。}
    \begin{itemize}
        \item CDS买方仍然面临风险：
        \begin{enumerate}
            \item 参考债务的信用风险
            \item CDS卖方的信用风险
            \item 基础风险（basis risk）
        \end{enumerate}
    \end{itemize}
    
    \textbf{B选项正确。}
    \begin{itemize}
        \item CDS买方确实面临卖方的信用风险，因为：
        \begin{enumerate}
            \item 如果参考债务违约，卖方可能无法支付所需款项
            \item 两个因素影响这种风险：
            \begin{itemize}
                \item 参考债务违约概率随时间增加，卖方违约概率也可能增加
                \item 经济和运营因素可能导致参考债务违约和卖方违约之间的相关性增加
            \end{itemize}
        \end{enumerate}
    \end{itemize}
    
    \textbf{C选项错误。}
    \begin{itemize}
        \item CDS买方不需要向卖方提供抵押品：
        \begin{enumerate}
            \item 买方支付固定费用（CDS spread）
            \item 抵押品要求通常由卖方提供
        \end{enumerate}
    \end{itemize}
    
    \textbf{D选项错误。}
    \begin{itemize}
        \item CDS买方不能以累积支付的价值放弃CDS：
        \begin{enumerate}
            \item CDS是双边合约
            \item 没有提前终止的自动机制
            \item 需要双方同意才能终止
        \end{enumerate}
    \end{itemize}
\end{explanation}

\checklist{CDS买方面临卖方的信用风险，因为如果参考债务违约，卖方可能无法支付所需款项}

\begin{question}
    With regard to the factors that determine recovery rates of traded bonds, which of the following statements is false?
\end{question}

\begin{choices}
    \item The liquidity of the debtor's collateral has an impact of recovery rates of the firm's traded bonds
    \item The presence of more bank loans increases the recovery rates of the firm's traded bonds
    \item The seniority of debtor's debt has an impact of recovery rates of the firm's traded bonds
    \item The value of the debtor's collateral has an impact of recovery rates of the firm's traded bonds
\end{choices}

\begin{correctanswer}
    B
\end{correctanswer}

\begin{explanation}
    \textbf{A选项错误。}
    \begin{itemize}
        \item 债务人抵押品的流动性确实会影响交易债券的回收率：
        \begin{enumerate}
            \item 流动性高的抵押品更容易变现
            \item 流动性低的抵押品可能难以在违约时快速出售
        \end{enumerate}
    \end{itemize}
    
    \textbf{B选项正确。}
    \begin{itemize}
        \item 更多的银行贷款实际上会降低交易债券的回收率，因为：
        \begin{enumerate}
            \item 银行贷款通常是优先级债务
            \item 银行贷款通常有最高级别的担保
            \item 公司使用的银行融资越多，交易债券的回收率就越低
        \end{enumerate}
    \end{itemize}
    
    \textbf{C选项错误。}
    \begin{itemize}
        \item 债务人债务的优先级确实会影响交易债券的回收率：
        \begin{enumerate}
            \item 优先级定义了公司违约时债权人对资产的优先权顺序
            \item 优先级高的债务通常有更高的回收率
        \end{enumerate}
    \end{itemize}
    
    \textbf{D选项错误。}
    \begin{itemize}
        \item 债务人抵押品的价值确实会影响交易债券的回收率：
        \begin{enumerate}
            \item 抵押品价值越高，回收率可能越高
            \item 抵押品价值是回收率的重要决定因素
        \end{enumerate}
    \end{itemize}
\end{explanation}

\checklist{更多的银行贷款会降低交易债券的回收率，因为银行贷款通常是优先级债务且有最高级别的担保}


\begin{question}
    A six-year CDS on an AA-rated issuer is offered at 150bp with semiannual payments while the yield on a six-year semiannual coupon bond of this issuer is 8\%. There is no counterparty risk on the CDS. The annualized LIBOR rate paid every six months is 4.6\% for all maturities. Which strategy would exploit the arbitrage opportunity? How much would your return exceed LIBOR?
\end{question}

\begin{choices}
    \item Buy the bond and the CDS with a risk-free gain of 1.9\%
    \item Buy the bond and the CDS with a risk-free gain of 0.32\%
    \item Short the bond and sell CDS protection with a risk-free gain of 4.97\%
    \item There is no arbitrage opportunity as any apparent risk-free profit is necessarily compensation for being exposed to the credit risk of the issuer
\end{choices}

\begin{correctanswer}
    A
\end{correctanswer}

\begin{explanation}
    \textbf{A选项正确。}
    \begin{itemize}
        \item 由于LIBOR持平，固定息票率为4.6\%，在债券上产生了340bp的价差：
        \begin{enumerate}
            \item 债券收益率 = 8\% = 800bp
            \item LIBOR = 4.6\% = 460bp
            \item 债券价差 = 800 - 460 = 340bp
        \end{enumerate}
        \item 通过做多债券和购买CDS可以产生190bp的年化利润：
        \begin{enumerate}
            \item CDS spread = 150bp
            \item 年化利润 = 340 - 150 = 190bp = 1.9\%
        \end{enumerate}
    \end{itemize}
    
    \textbf{B选项错误。}
    \begin{itemize}
        \item 0.32\%的收益计算错误
        \item 正确的收益应该是1.9\%
    \end{itemize}
    
    \textbf{C选项错误。}
    \begin{itemize}
        \item 做空债券和出售CDS保护会产生损失
        \item 4.97\%的收益计算错误
    \end{itemize}
    
    \textbf{D选项错误。}
    \begin{itemize}
        \item 确实存在套利机会
        \item 由于CDS没有交易对手风险，这个套利是真正无风险的
    \end{itemize}
\end{explanation}

\checklist{通过做多债券和购买CDS可以产生1.9\%的无风险收益，因为债券价差（340bp）大于CDS spread（150bp）}

\begin{question}
    The Merton model and the Moody's KMV model use different approaches to determine the probability of default. Which of the following is consistent with Moody's KMV model?
\end{question}

\begin{choices}
    \item The distance to default is 1.96, so there is a 2.5\% probability of default
    \item The distance to default is 1.96, so there is a 5.0\% probability of default
    \item The historical frequency of default for corporate bonds has been 6\%. Updating this with Altman's Z-score analysis would provide a probability of default that is somewhat different than 6\%
    \item The distance to default is 1.96 and, historically, 1.2\% of firms with this characterization have defaulted, so there is a 1.2\% probability of default
\end{choices}

\begin{correctanswer}
    D
\end{correctanswer}

\begin{explanation}
    \textbf{A选项错误。}
    \begin{itemize}
        \item 2.5\%的违约概率是基于正态分布假设
        \item Moody's KMV模型不使用正态分布假设
        \item 它使用历史违约数据来确定违约概率
    \end{itemize}
    
    \textbf{B选项错误。}
    \begin{itemize}
        \item 5.0\%的违约概率也是基于正态分布假设
        \item 这与Moody's KMV模型的方法不符
        \item 模型使用历史违约频率而不是理论概率
    \end{itemize}
    
    \textbf{C选项错误。}
    \begin{itemize}
        \item 使用Altman's Z-score分析不是Moody's KMV模型的方法
        \item 模型使用距离违约（distance to default）作为主要指标
        \item 不依赖于Z-score分析
    \end{itemize}
    
    \textbf{D选项正确。}
    \begin{itemize}
        \item Moody's KMV模型：
        \begin{enumerate}
            \item 评估具有相似违约距离的公司的历史违约频率
            \item 使用这个历史频率作为违约概率
            \item 不依赖于理论分布假设
        \end{enumerate}
        \item 在这个例子中：
        \begin{enumerate}
            \item 违约距离为1.96
            \item 历史上具有这个特征的公司有1.2\%违约
            \item 因此违约概率为1.2\%
        \end{enumerate}
    \end{itemize}
\end{explanation}

\checklist{Moody's KMV模型使用具有相似违约距离的公司的历史违约频率来确定违约概率，而不是使用理论分布假设}


\begin{question}
    A corporate bond with a face value of \$1,000 has a remaining maturity of 10 years. Using the Merton model, the current value of the bond is calculated at \$750. Assuming that the risk-free rate is equal to 3\%, what is the credit spread for this bond?
\end{question}

\begin{choices}
    \item 35 bps
    \item 42 bps
    \item 58 bps
    \item 65 bps
\end{choices}

\begin{correctanswer}
    C
\end{correctanswer}

\begin{explanation}
    \textbf{A选项错误。}
    \begin{itemize}
        \item 35个基点的计算错误
        \item 没有正确使用Merton模型的信用利差公式
    \end{itemize}
    
    \textbf{B选项错误。}
    \begin{itemize}
        \item 42个基点的计算错误
        \item 可能混淆了债券现值（\$750）和信用利差
    \end{itemize}
    
    \textbf{C选项正确。}
    \begin{itemize}
        \item 使用Merton模型的信用利差公式：
        \[
        \begin{aligned}
        \text{信用利差} &= -\left(\frac{1}{T-t}\right) \ln\left(\frac{D}{F}\right) - r \\
        &= -\left(\frac{1}{10}\right) \times \ln\left(\frac{750}{1000}\right) - 0.03 \\
        &= 58\ \text{个基点}
        \end{aligned}
        \]
        \item 其中：
        \begin{enumerate}
            \item $T-t = 10$年（剩余期限）
            \item $D = \$750$（债券现值）
            \item $F = \$1,000$（面值）
            \item $r = 3\%$（无风险利率）
        \end{enumerate}
    \end{itemize}
    
    \textbf{D选项错误。}
    \begin{itemize}
        \item 65个基点的计算错误
        \item 可能使用了错误的公式或参数
    \end{itemize}
\end{explanation}

\checklist{使用Merton模型的信用利差公式，通过债券现值、面值、剩余期限和无风险利率计算得出58个基点的信用利差}

\begin{question}
    The capital structure of ABC Corporation consists of two parts, one 4 year bond with face value of \$120 million and the rest is equity. The current market value of the firm assets is \$150 and the expected rate of change of the firm's value is 20\%. The volatility is 25\%. The firm's risk management division estimates the distance to default using Merton model. Given the distance to default, the estimated physical default probability is? (N(1.85) = 96.78\%; N(2.45) = 99.29\%)
\end{question}

\begin{choices}
    \item 3.22\%
    \item 15.36\%
    \item 15.38\%
    \item 28.45\%
\end{choices}

\begin{correctanswer}
    A
\end{correctanswer}

\begin{explanation}
    \textbf{A选项正确。}
    \begin{itemize}
        \item 使用Merton模型计算违约距离：
        \[
        \begin{aligned}
        d_2 &= \frac{\ln\left[\frac{150}{120} \times e^{-20\% \times 4}\right]}{0.25 \times \sqrt{4}} - \frac{1}{2} \times 0.25 \times \sqrt{4} \\
        &= 1.85
        \end{aligned}
        \]
        \item 实际违约概率 = N(-d2) = 3.22\%
        \item 其中：
        \begin{enumerate}
            \item 公司资产价值 = \$150 million
            \item 债券面值 = \$120 million
            \item 预期增长率 = 20\%
            \item 波动率 = 25\%
            \item 期限 = 4年
        \end{enumerate}
    \end{itemize}
    
    \textbf{B选项错误。}
    \begin{itemize}
        \item 15.36\%的计算错误
        \item 可能使用了错误的公式或参数
    \end{itemize}
    
    \textbf{C选项错误。}
    \begin{itemize}
        \item 15.38\%的计算错误
        \item 可能混淆了违约距离和违约概率
    \end{itemize}
    
    \textbf{D选项错误。}
    \begin{itemize}
        \item 28.45\%的计算错误
        \item 可能直接使用了波动率作为违约概率
    \end{itemize}
\end{explanation}

\checklist{使用Merton模型计算违约距离，通过公司资产价值、债券面值、预期增长率、波动率和期限计算得出3.22\%的实际违约概率}


\begin{question}
    Which of the following statements regarding the Merton model is true?
\end{question}

\begin{choices}
    \item A firm with multiple debt issues that mature at different times is easy to value with the Merton model
    \item The Merton model assumes a lognormal distribution and constant variance for changes in firm value
    \item The Merton model is able to predict default because it allows for default surprises (i.e., jumps)
    \item Empirical results indicate that the Merton model is able to predict default better than naive models for investment grade bonds
\end{choices}

\begin{correctanswer}
    B
\end{correctanswer}

\begin{explanation}
    \textbf{A选项错误。}
    \begin{itemize}
        \item Merton模型难以处理多种债务工具：
        \begin{enumerate}
            \item 大多数公司有不同到期日的债务
            \item 债务可能有不同的票面利率
            \item 模型假设单一零息债券，这与现实不符
        \end{enumerate}
    \end{itemize}
    
    \textbf{B选项正确。}
    \begin{itemize}
        \item Merton模型的基本假设：
        \begin{enumerate}
            \item 公司价值服从对数正态分布
            \item 公司价值变化具有恒定方差
            \item 这些假设是模型的核心特征
        \end{enumerate}
    \end{itemize}
    
    \textbf{C选项错误。}
    \begin{itemize}
        \item Merton模型不允许公司价值跳跃：
        \begin{enumerate}
            \item 模型假设公司价值连续变化
            \item 不能处理突然的违约事件
            \item 这使得违约预测过于可预测
        \end{enumerate}
    \end{itemize}
    
    \textbf{D选项错误。}
    \begin{itemize}
        \item 实证研究结果：
        \begin{enumerate}
            \item 对于投资级债券，简单模型（假设债务无风险）比Merton模型更有效
            \item Merton模型对非投资级债券的违约预测更有效
            \item 这与选项描述相反
        \end{enumerate}
    \end{itemize}
\end{explanation}

\checklist{Merton模型假设公司价值服从对数正态分布且具有恒定方差，这是模型的核心特征}


\begin{question}
    Which of the following statements about credit default swaps is most accurate?
\end{question}

\begin{choices}
    \item CDSs transfer credit risk and market risk from the protection buyer to the protection seller
    \item CDSs transfer credit risk from the protection buyer to the issuer of the underlying credit
    \item Physical settlement requires knowledge of the post-default market price
    \item Cash settlement avoids the problem of a delivery squeeze
\end{choices}

\begin{correctanswer}
    D
\end{correctanswer}

\begin{explanation}
    \textbf{A选项错误。}
    \begin{itemize}
        \item CDS只转移信用风险：
        \begin{enumerate}
            \item 不转移市场风险（如利率风险、汇率风险等）
            \item 保护买方仍然面临市场风险
            \item 保护卖方只承担信用风险
        \end{enumerate}
    \end{itemize}
    
    \textbf{B选项错误。}
    \begin{itemize}
        \item CDS将信用风险转移给保护卖方：
        \begin{enumerate}
            \item 不是转移给基础信用的发行人
            \item 发行人仍然是信用风险的源头
            \item 保护卖方承担信用风险
        \end{enumerate}
    \end{itemize}
    
    \textbf{C选项错误。}
    \begin{itemize}
        \item 实物结算不需要知道违约后的市场价格：
        \begin{enumerate}
            \item 基于参考债务的面值
            \item 不需要市场价格信息
            \item 这是实物结算的优势之一
        \end{enumerate}
    \end{itemize}
    
    \textbf{D选项正确。}
    \begin{itemize}
        \item 现金结算的优势：
        \begin{enumerate}
            \item 没有实际的证券交易
            \item 避免了交割挤压问题
            \item 保护买方不需要在公开市场购买参考实体
            \item 不会因为大量购买导致价格上升
        \end{enumerate}
    \end{itemize}
\end{explanation}

\checklist{现金结算避免了交割挤压问题，因为不需要实际的证券交易，保护买方不需要在公开市场购买参考实体}


\begin{question}
    A 3-year credit default swap (CDS) specifying physical delivery defaults at the end of three years. If the reference asset is a \$250 million, 7.5\% ABC corporate bond, and the CDS spread is 150 basis points, the buyer of the CDS will:
\end{question}

\begin{choices}
    \item Receive payments of 750 basis points for the next three years
    \item Receive a payment of \$187.5 million
    \item Deliver the bond and receive a payment of \$250 million
    \item Continue to receive payments of 600 basis points for the next three years
\end{choices}

\begin{correctanswer}
    C
\end{correctanswer}

\begin{explanation}
    \textbf{A选项错误。}
    \begin{itemize}
        \item CDS买方不是接收支付，而是支付固定费用：
        \begin{enumerate}
            \item 750个基点是债券的票面利率
            \item 买方需要支付CDS spread（150个基点）给卖方
        \end{enumerate}
    \end{itemize}
    
    \textbf{B选项错误。}
    \begin{itemize}
        \item \$187.5 million不是正确的支付金额：
        \begin{enumerate}
            \item 在实物交割的CDS中，买方会收到参考债务的面值
            \item 面值是\$250 million
        \end{enumerate}
    \end{itemize}
    
    \textbf{C选项正确。}
    \begin{itemize}
        \item 在实物交割的CDS中：
        \begin{enumerate}
            \item CDS买方将参考债务（ABC公司债券）交付给卖方
            \item CDS买方收到参考债务的面值（\$250 million）
            \item 这是实物交割CDS的标准结算方式
        \end{enumerate}
    \end{itemize}
    
    \textbf{D选项错误。}
    \begin{itemize}
        \item CDS买方不是接收支付：
        \begin{enumerate}
            \item 买方需要支付CDS spread（150个基点）给卖方
            \item 600个基点不是正确的CDS spread
        \end{enumerate}
    \end{itemize}
\end{explanation}

\checklist{实物交割的CDS中，买方将参考债务交付给卖方，并收到参考债务的面值（\$250 million）}


\begin{question}
    A default swap acts like a:
\end{question}

\begin{choices}
    \item Call option on the reference obligation for the buyer of the swap
    \item Put option on the reference obligation for the buyer of the swap
    \item Look back option on the reference obligation for the buyer of the swap
    \item Up-and-out option on the reference obligation for the buyer of the swap
\end{choices}

\begin{correctanswer}
    B
\end{correctanswer}

\begin{explanation}
    \textbf{A选项错误。}
    \begin{itemize}
        \item 看涨期权（Call Option）的特征：
        \begin{enumerate}
            \item 买方有权以固定价格买入资产
            \item 当资产价格上升时获利
            \item 这与违约互换的功能不符
        \end{enumerate}
    \end{itemize}
    
    \textbf{B选项正确。}
    \begin{itemize}
        \item 违约互换类似于看跌期权（Put Option）：
        \begin{enumerate}
            \item 买方在违约时获得支付
            \item 限制了买方的下行风险
            \item 类似于看跌期权的保护功能
        \end{enumerate}
        \item 具体表现：
        \begin{enumerate}
            \item 当参考债务价值下降时提供保护
            \item 在违约事件发生时获得补偿
            \item 类似于看跌期权的行权机制
        \end{enumerate}
    \end{itemize}
    
    \textbf{C选项错误。}
    \begin{itemize}
        \item 回顾期权（Look Back Option）的特征：
        \begin{enumerate}
            \item 基于资产价格的历史表现
            \item 与违约事件无关
            \item 不适用于违约互换
        \end{enumerate}
    \end{itemize}
    
    \textbf{D选项错误。}
    \begin{itemize}
        \item 向上敲出期权（Up-and-out Option）的特征：
        \begin{enumerate}
            \item 当资产价格超过特定水平时失效
            \item 与违约事件无关
            \item 不适用于违约互换
        \end{enumerate}
    \end{itemize}
\end{explanation}

\checklist{违约互换的行为类似于看跌期权，买方在违约时获得支付，限制了买方的下行风险}


\begin{question}
    Which of the following statements is most likely correct about recovery rates?
\end{question}

\begin{choices}
    \item Recovery rates vary inversely with capital structure seniority
    \item Historical recovery rates are fairly constant across industries
    \item Recovery rates are highest during economic downturns
    \item Actual recovery rates may differ substantially from settled recovery rates
\end{choices}

\begin{correctanswer}
    D
\end{correctanswer}

\begin{explanation}
    \textbf{A选项错误。}
    \begin{itemize}
        \item 回收率与资本结构优先级的关系：
        \begin{enumerate}
            \item 回收率随优先级的增加而增加
            \item 优先级高的债务通常有更高的回收率
            \item 不是反向关系
        \end{enumerate}
    \end{itemize}
    
    \textbf{B选项错误。}
    \begin{itemize}
        \item 不同行业的回收率差异：
        \begin{enumerate}
            \item 行业间回收率差异显著
            \item 资产密集型行业通常有更高的回收率
            \item 服务型行业通常有较低的回收率
        \end{enumerate}
    \end{itemize}
    
    \textbf{C选项错误。}
    \begin{itemize}
        \item 经济衰退期间的回收率：
        \begin{enumerate}
            \item 经济衰退期间回收率通常最低
            \item 资产价值下降导致回收率降低
            \item 市场流动性不足影响回收率
        \end{enumerate}
    \end{itemize}
    
    \textbf{D选项正确。}
    \begin{itemize}
        \item 结算回收率与实际回收率的差异：
        \begin{enumerate}
            \item 结算回收率在信用事件后很快发生（如CDS拍卖或违约债券出售）
            \item 实际回收率可能发生在多年后（基于破产解决）
            \item 两种回收率可能存在显著差异
        \end{enumerate}
    \end{itemize}
\end{explanation}

\checklist{实际回收率可能与结算回收率存在显著差异，因为结算回收率在信用事件后很快发生，而实际回收率可能发生在多年后}

\begin{question}
    Which of the following statements is most likely correct about recovery rates?
\end{question}

\begin{choices}
    \item Recovery rates are higher for secured debt than for unsecured debt
    \item Recovery rates remain stable across different economic cycles
    \item Recovery rates are higher for subordinated debt than for senior debt
    \item Actual recovery rates may differ substantially from settled recovery rates
\end{choices}

\begin{correctanswer}
    D
\end{correctanswer}

\begin{explanation}
    \textbf{A选项错误。}
    \begin{itemize}
        \item 有担保债务的回收率：
        \begin{enumerate}
            \item 有担保债务确实通常有更高的回收率
            \item 但这只是影响回收率的一个因素
            \item 不是最准确的描述
        \end{enumerate}
    \end{itemize}
    
    \textbf{B选项错误。}
    \begin{itemize}
        \item 经济周期对回收率的影响：
        \begin{enumerate}
            \item 回收率随经济周期变化
            \item 经济衰退期间回收率通常较低
            \item 经济繁荣期间回收率通常较高
        \end{enumerate}
    \end{itemize}
    
    \textbf{C选项错误。}
    \begin{itemize}
        \item 次级债务与优先债务的回收率：
        \begin{enumerate}
            \item 优先债务通常有更高的回收率
            \item 次级债务的回收率通常较低
            \item 与选项描述相反
        \end{enumerate}
    \end{itemize}
    
    \textbf{D选项正确。}
    \begin{itemize}
        \item 结算回收率与实际回收率的差异：
        \begin{enumerate}
            \item 结算回收率在信用事件后很快发生（如CDS拍卖或违约债券出售）
            \item 实际回收率可能发生在多年后（基于破产解决）
            \item 两种回收率可能存在显著差异
        \end{enumerate}
    \end{itemize}
\end{explanation}

\checklist{实际回收率可能与结算回收率存在显著差异，因为结算回收率在信用事件后很快发生，而实际回收率可能发生在多年后}

\begin{question}
    Suppose a firm has two debt issues outstanding. One is a senior debt issue that matures in four years with a principal amount of \$120 million. The other is a subordinate debt issue that also matures in four years with a principal amount of \$60 million. The annual interest rate is 4\% and the volatility of the firm value is estimated to be 20\%. If the volatility of the firm value declines in the Merton model then which of the following statements is true?
\end{question}

\begin{choices}
    \item If the firm is experiencing financial distress (low firm value), then the value of senior debt will increase while the values of subordinate debt and equity will both decline
    \item If the firm is not experiencing financial distress (high firm value), then the value of senior debt and subordinate debt and equity will increase
    \item If the firm is experiencing financial distress (low firm value), then the value of senior debt and subordinate debt will increase while equity values will decline
    \item If the firm is not experiencing financial distress (high firm value), then the value of senior debt will increase while the values of subordinate debt and equity will both decline
\end{choices}

\begin{correctanswer}
    A
\end{correctanswer}

\begin{explanation}
    \textbf{A选项正确。}
    \begin{itemize}
        \item 当公司经历财务困境时：
        \begin{enumerate}
            \item 波动率下降会降低股权价值（因为股权被视为看涨期权）
            \item 波动率下降会增加高级债价值
            \item 次级债的表现类似股权，其价值也会下降
        \end{enumerate}
        \item 原因：
        \begin{enumerate}
            \item 在财务困境时，次级债的风险特征更接近股权
            \item 波动率下降减少了期权的价值
            \item 高级债作为剩余索取权，价值上升
        \end{enumerate}
    \end{itemize}
    
    \textbf{B选项错误。}
    \begin{itemize}
        \item 当公司没有财务困境时：
        \begin{enumerate}
            \item 波动率下降会增加高级债价值
            \item 次级债的表现类似高级债，其价值也会上升
            \item 股权价值会下降
        \end{enumerate}
    \end{itemize}
    
    \textbf{C选项错误。}
    \begin{itemize}
        \item 在财务困境时：
        \begin{enumerate}
            \item 次级债价值会下降，而不是上升
            \item 因为次级债的表现类似股权
        \end{enumerate}
    \end{itemize}
    
    \textbf{D选项错误。}
    \begin{itemize}
        \item 在没有财务困境时：
        \begin{enumerate}
            \item 次级债价值会上升，而不是下降
            \item 因为次级债的表现类似高级债
        \end{enumerate}
    \end{itemize}
\end{explanation}

\checklist{在财务困境时，波动率下降会导致高级债价值上升，而次级债和股权价值下降，因为次级债的风险特征更接近股权}

\begin{question}
    Which of the following statements regarding the Merton model is true?
\end{question}

\begin{choices}
    \item A firm with multiple debt issues that mature at different times is easy to value with the Merton model
    \item The Merton model assumes a lognormal distribution and constant variance for changes in firm value
    \item The Merton model is able to predict default because it allows for default surprises (i.e., jumps)
    \item Empirical results indicate that the Merton model is able to predict default better than naive models for investment grade bonds
\end{choices}

\begin{correctanswer}
    B
\end{correctanswer}

\begin{explanation}
    \textbf{A选项错误。}
    \begin{itemize}
        \item Merton模型的局限性：
        \begin{enumerate}
            \item 假设单一零息债券
            \item 难以处理多种债务工具
            \item 难以处理不同到期日的债务
            \item 难以处理不同票面利率的债务
        \end{enumerate}
    \end{itemize}
    
    \textbf{B选项正确。}
    \begin{itemize}
        \item Merton模型的基本假设：
        \begin{enumerate}
            \item 公司价值服从对数正态分布
            \item 公司价值变化具有恒定方差
            \item 这些假设是模型的核心特征
        \end{enumerate}
    \end{itemize}
    
    \textbf{C选项错误。}
    \begin{itemize}
        \item Merton模型的违约预测：
        \begin{enumerate}
            \item 不允许公司价值跳跃
            \item 假设公司价值连续变化
            \item 不能处理突然的违约事件
            \item 这使得违约预测过于可预测
        \end{enumerate}
    \end{itemize}
    
    \textbf{D选项错误。}
    \begin{itemize}
        \item 实证研究结果：
        \begin{enumerate}
            \item 对于投资级债券，简单模型（假设债务无风险）比Merton模型更有效
            \item Merton模型对非投资级债券的违约预测更有效
            \item 这与选项描述相反
        \end{enumerate}
    \end{itemize}
\end{explanation}

\checklist{Merton模型假设公司价值服从对数正态分布且具有恒定方差，这是模型的核心特征}


\begin{question}
    Each of the following items represents an example of qualitative information that would ideally be captured in assessing default probability except:
\end{question}

\begin{choices}
    \item Management's education and experience
    \item Internal controls associated with financial reporting
    \item Diversification of products and customers locally and globally
    \item Trends in throughput and other operational efficiency metrics
\end{choices}

\begin{correctanswer}
    D
\end{correctanswer}

\begin{explanation}
    \textbf{A选项错误。}
    \begin{itemize}
        \item 管理层的教育和经验是定性信息：
        \begin{enumerate}
            \item 难以用数字准确衡量
            \item 需要主观评估
            \item 对违约概率有重要影响
        \end{enumerate}
    \end{itemize}
    
    \textbf{B选项错误。}
    \begin{itemize}
        \item 财务报告相关的内部控制是定性信息：
        \begin{enumerate}
            \item 难以量化
            \item 需要专业判断
            \item 影响财务信息的可靠性
        \end{enumerate}
    \end{itemize}
    
    \textbf{C选项错误。}
    \begin{itemize}
        \item 产品和客户的多样化是定性信息：
        \begin{enumerate}
            \item 难以用单一指标衡量
            \item 需要综合评估
            \item 影响公司的风险分散程度
        \end{enumerate}
    \end{itemize}
    
    \textbf{D选项正确。}
    \begin{itemize}
        \item 吞吐量和其他运营效率指标是定量信息：
        \begin{enumerate}
            \item 可以用具体数字衡量
            \item 可以定期跟踪和比较
            \item 可以转化为标准化的指标
        \end{enumerate}
        \item 这些指标：
        \begin{enumerate}
            \item 可以量化
            \item 可以客观衡量
            \item 可以用于统计分析
        \end{enumerate}
    \end{itemize}
\end{explanation}

\checklist{吞吐量和其他运营效率指标是定量信息，可以用具体数字衡量，而其他选项都是需要主观评估的定性信息}


\begin{question}
    The rate parameter in the exponential distribution measures the rate at which it takes an event to occur. In the context of waiting for a company to default, the rate parameter is known as the hazard rate and indicates the rate at which default will arrive. Which of the following statements about hazard rates (i.e., default intensity) is correct assuming a constant default intensity of 0.15?
\end{question}

\begin{choices}
    \item The cumulative default probability after 3 periods = e-(0.15)3
    \item The conditional default probability after 1 period = 15\%
    \item The unconditional default probability is memoryless
    \item The conditional default probability is memoryless
\end{choices}

\begin{correctanswer}
    D
\end{correctanswer}

\begin{explanation}
    \textbf{A选项错误。}
    \begin{itemize}
        \item 累积违约概率的计算：
        \begin{enumerate}
            \item 正确公式：1 - e-(0.15)3
            \item 选项给出的是生存概率公式
            \item 违约强度为0.15时的累积违约概率约为36.1\%
        \end{enumerate}
    \end{itemize}
    
    \textbf{B选项错误。}
    \begin{itemize}
        \item 第一期的条件违约概率：
        \begin{enumerate}
            \item 第一期无条件违约概率 = 1 - e-(0.15)1 = 13.9\%
            \item 第二期累积违约概率 = 1 - e-(0.15)2 = 25.9\%
            \item 第二期条件违约概率 = (25.9\% - 13.9\%)/(1 - 13.9\%) = 14.0\%
        \end{enumerate}
    \end{itemize}
    
    \textbf{C选项错误。}
    \begin{itemize}
        \item 无条件违约概率：
        \begin{enumerate}
            \item 不是无记忆的
            \item 随时间累积
            \item 受历史违约情况影响
        \end{enumerate}
    \end{itemize}
    
    \textbf{D选项正确。}
    \begin{itemize}
        \item 条件违约概率：
        \begin{enumerate}
            \item 是无记忆的
            \item 只取决于当前状态
            \item 不受历史违约情况影响
        \end{enumerate}
        \item 原因：
        \begin{enumerate}
            \item 指数分布的无记忆性
            \item 违约强度恒定
            \item 未来违约概率独立于过去
        \end{enumerate}
    \end{itemize}
\end{explanation}

\checklist{条件违约概率是无记忆的，只取决于当前状态，不受历史违约情况影响，这是指数分布的特性}


\begin{question}
    Using the Merton model, calculate the current value of a firm's equity and debt given that the current value of the firm is \$120 million, the principal amount due in four years on the zero-coupon bond is \$100 million, the annual interest rate is 8\%, and the volatility of the firm is 25\%.
\end{question}

\begin{choices}
    \item \$100 million in debt and \$20 million in equity
    \item \$65.32 million in debt and \$54.68 million in equity
    \item \$62.45 million in debt and \$57.55 million in equity
    \item \$45.28 million in debt and \$74.72 million in equity
\end{choices}

\begin{correctanswer}
    C
\end{correctanswer}

\begin{explanation}
    \textbf{A选项错误。}
    \begin{itemize}
        \item 简单用公司价值减去债务面值：
        \begin{enumerate}
            \item 没有考虑债务的现值
            \item 没有考虑波动率的影响
            \item 没有使用Merton模型
        \end{enumerate}
    \end{itemize}
    
    \textbf{B选项错误。}
    \begin{itemize}
        \item 计算错误：
        \begin{enumerate}
            \item 可能使用了错误的贴现率
            \item 可能使用了错误的波动率
            \item 可能使用了错误的期限
        \end{enumerate}
    \end{itemize}
    
    \textbf{C选项正确。}
    \begin{itemize}
        \item 使用Merton模型计算：
        \[
        \begin{aligned}
        S_t &= 120N(d) - (100)(0.7350)N(d - \sigma\sqrt{T-t}) \\
        d &= \frac{\ln\left(\frac{120}{100 \times 0.7350}\right)}{0.25\sqrt{4}} + \frac{1}{2}(0.25)\sqrt{4} \\
        &= \frac{\ln(1.6327)}{0.5} + \frac{1}{2}(0.5) = 1.2456
        \end{aligned}
        \]
        \[
        \begin{aligned}
        S_t &= 120N(1.2456) - (73.50)N(1.2456 - 0.5) \\
        &= 120(0.8934) - (73.50)(0.6915) \\
        &= 107.21 - 50.83 = 57.55
        \end{aligned}
        \]
        \item 因此：
        \begin{enumerate}
            \item 股权价值 = \$57.55 million
            \item 债务价值 = \$120 - \$57.55 = \$62.45 million
        \end{enumerate}
    \end{itemize}
    
    \textbf{D选项错误。}
    \begin{itemize}
        \item 计算错误：
        \begin{enumerate}
            \item 可能低估了债务价值
            \item 可能高估了股权价值
            \item 可能使用了错误的参数
        \end{enumerate}
    \end{itemize}
\end{explanation}

\checklist{使用Merton模型计算得出股权价值、债务价值，考虑公司价值、债务面值、利率、波动率和期限的影响}

\begin{question}
    A financial analyst wants to estimate the value of a company's equity using the Merton model. The company pays no dividends and has a market value of assets equal to EUR 80 million. Its only debt, a zero-coupon bond maturing in 3 years, has a face value of EUR 50 million and a current price of EUR 62.50 per EUR 100 of face value. Additional information is given below:
    • N(d1): 0.9625
    • N(d2): 0.9012
    • Annual risk-free interest rate: 4\%
    What is the value of the company's equity?
\end{question}

\begin{choices}
    \item EUR 28.6 million
    \item EUR 32.4 million
    \item EUR 42.8 million
    \item EUR 47.5 million
\end{choices}

\begin{correctanswer}
    D
\end{correctanswer}

\begin{explanation}
    \textbf{A选项错误。}
    \begin{itemize}
        \item 计算错误：
        \begin{enumerate}
            \item 可能使用了错误的N(d1)和N(d2)值
            \item 可能没有正确考虑债务现值
            \item 可能使用了错误的公司价值
        \end{enumerate}
    \end{itemize}
    
    \textbf{B选项错误。}
    \begin{itemize}
        \item 计算错误：
        \begin{enumerate}
            \item 可能低估了股权价值
            \item 可能高估了债务价值
            \item 可能使用了错误的贴现率
        \end{enumerate}
    \end{itemize}
    
    \textbf{C选项错误。}
    \begin{itemize}
        \item 计算错误：
        \begin{enumerate}
            \item 可能使用了错误的债务面值
            \item 可能没有正确考虑债务价格
            \item 可能使用了错误的期限
        \end{enumerate}
    \end{itemize}
    
    \textbf{D选项正确。}
    \begin{itemize}
        \item 使用Merton模型计算：
        \[
        \begin{aligned}
        S_t &= V_tN(d_1) - Ke^{-r(T-t)}N(d_2) \\
        &= 80 \times 0.9625 - 50 \times e^{-0.04 \times 3} \times 0.9012 \\
        &= 77 - 29.5 = 47.5
        \end{aligned}
        \]
        \item 其中：
        \begin{enumerate}
            \item 公司价值 = EUR 80 million
            \item 债务面值 = EUR 50 million
            \item 债务现值 = EUR 31.25 million (50 × 62.50\%)
            \item 无风险利率 = 4\%
            \item 期限 = 3年
        \end{enumerate}
    \end{itemize}
\end{explanation}

\checklist{使用Merton模型计算得出股权价值为EUR 47.5 million，其中考虑了公司价值、债务面值、债务价格、无风险利率和期限的影响}

\begin{question}
    The CRO of Bank KMZ has asked the bank's risk managers to assess intraday liquidity risk and estimate capital requirements for the bank under the assumption that stressed market conditions will occur over the next 1 year. The CRO instructs the managers to model a parallel shift of 280 bps in the yield curve even though nonparallel shifts in the yield curve are plausible. Bank KMZ holds a portfolio of credit assets of the same type and applies a multiplier factor for both VaR and stressed VaR (SVaR) equal to 3.0. The estimated VaR, SVaR, and the average VaR factors for the bank, obtained with a 10-day time horizon and 99% confidence level, are reported in the table below, along with 1-year default rate at the 99.9% level, probability of default (PD), exposure at default (EAD), and recovery rate data:

    Which of the following statements is correct? CVaR = WCL - EL = 105 × 0.95
\end{question}

\begin{choices}
    \item The estimated credit risk capital requirement over a 1-year time horizon for Bank KMZ according to the Basel II Internal Ratings-Based Approach is EUR 0.8 million
    \item The estimated market risk capital requirement over 1-year time horizon for Bank KMZ according to Basel 2.5 is EUR 24.5 million
    \item The Basel guidance on stress testing of intraday liquidity risk management requires Bank KMZ to model only two scenarios: market-wide liquidity stress and the bank's own stress
    \item The total risk budget for Bank KMZ's asset portfolio will be the same whether Bank KMZ use VaR or SVaR as the basis for its budgeting decision
\end{choices}

\begin{correctanswer}
    B
\end{correctanswer}

\begin{explanation}
    \textbf{A选项错误。}
    \begin{itemize}
        \item 信用风险资本要求计算错误：
        \begin{enumerate}
            \item 没有考虑违约概率和违约损失
            \item 没有考虑风险敞口
            \item 没有考虑回收率
        \end{enumerate}
    \end{itemize}
    
    \textbf{B选项正确。}
    \begin{itemize}
        \item 市场风险资本要求计算：
        \begin{enumerate}
            \item 使用VaR和SVaR的较大值
            \item 考虑乘数因子3.0
        \end{enumerate}
        \item 计算过程：
        \begin{enumerate}
            \item 10天VaR = 2.3 million
            \item 10天SVaR = 3.5 million
            \item 年化VaR = 2.3 × √25 = 11.5 million
            \item 年化SVaR = 3.5 × √25 = 17.5 million
            \item 资本要求 = 17.5 × 3.0 = 24.5 million
        \end{enumerate}
    \end{itemize}
    
    \textbf{C选项错误。}
    \begin{itemize}
        \item 巴塞尔委员会要求：
        \begin{enumerate}
            \item 需要模拟多个情景
            \item 不仅限于两个情景
            \item 包括系统性风险和特定风险
        \end{enumerate}
    \end{itemize}
    
    \textbf{D选项错误。}
    \begin{itemize}
        \item VaR和SVaR的区别：
        \begin{enumerate}
            \item SVaR通常大于VaR
            \item 风险预算会不同
            \item 使用SVaR会导致更高的资本要求
        \end{enumerate}
    \end{itemize}
\end{explanation}

\checklist{市场风险资本要求使用VaR和SVaR的较大值计算}


\begin{question}
    A large regional bank is changing the method it uses to calculate credit risk capital from the standardized approach to the Foundation Internal Ratings-Based (IRB) approach. A risk analyst is asked to summarize the process of calculating the capital charge under the new approach. Which of the following actions would the bank be required to take as part of the Foundation IRB approach?
\end{question}

\begin{choices}
    \item Apply an adjustment factor that reduces the risk weight for credit exposures with maturities longer than 1 year
    \item Develop a matrix to estimate the correlations between different classes of credit exposures
    \item Assume that correlations between high-risk corporate bonds are constant across all market conditions
    \item Provide an internally modeled default probability for each credit exposure
\end{choices}

\begin{correctanswer}
    D
\end{correctanswer}

\begin{explanation}
    \textbf{A选项错误。}
    \begin{itemize}
        \item 期限调整因子：
        \begin{enumerate}
            \item 不是基础IRB方法的要求
            \item 是高级IRB方法的特征
            \item 基础IRB方法使用标准期限调整
        \end{enumerate}
    \end{itemize}
    
    \textbf{B选项错误。}
    \begin{itemize}
        \item 相关性矩阵：
        \begin{enumerate}
            \item 不是基础IRB方法的要求
            \item 是高级IRB方法的特征
            \item 基础IRB方法使用标准相关性
        \end{enumerate}
    \end{itemize}
    
    \textbf{C选项错误。}
    \begin{itemize}
        \item 相关性假设：
        \begin{enumerate}
            \item 不是基础IRB方法的要求
            \item 相关性会随市场条件变化
            \item 基础IRB方法使用标准相关性
        \end{enumerate}
    \end{itemize}
    
    \textbf{D选项正确。}
    \begin{itemize}
        \item 基础IRB方法的要求：
        \begin{enumerate}
            \item 必须为每个信用敞口提供内部违约概率模型
            \item 这是基础IRB方法的核心要求
            \item 其他参数使用标准值
        \end{enumerate}
    \end{itemize}
\end{explanation}

\checklist{基础IRB方法要求为每个信用敞口提供内部违约概率模型}


\begin{question}
    A credit analyst is estimating a bank's economic capital for credit risk at the 99.9% confidence level. The bank reports that its current credit risk portfolio has the following parameters:
    · Total exposure amount: GBP 250 million
    · Probability of default (PD): 2.5\%
    · Loss rate (LR): 12\%
    · Standard deviation of the PD: 8\%
    · Standard deviation of the LR: 15\%
    · Capital multiplier for the portfolio at the 99.9% confidence level: 7.50
    What amount of economic capital should the bank hold for its credit risk portfolio?
\end{question}

\begin{choices}
    \item GBP 5.25 million
    \item GBP 6.75 million
    \item GBP 37.50 million
    \item GBP 45.00 million
\end{choices}

\begin{correctanswer}
    D
\end{correctanswer}

\begin{explanation}
    \textbf{A选项错误。}
    \begin{itemize}
        \item 计算错误：
        \begin{enumerate}
            \item 可能只考虑了预期损失
            \item 没有考虑非预期损失
            \item 没有使用资本乘数
        \end{enumerate}
    \end{itemize}
    
    \textbf{B选项错误。}
    \begin{itemize}
        \item 计算错误：
        \begin{enumerate}
            \item 可能使用了错误的乘数
            \item 可能低估了风险
            \item 没有考虑所有风险因素
        \end{enumerate}
    \end{itemize}
    
    \textbf{C选项错误。}
    \begin{itemize}
        \item 计算错误：
        \begin{enumerate}
            \item 可能只考虑了PD的标准差
            \item 可能忽略了LR的标准差
            \item 可能使用了错误的敞口金额
        \end{enumerate}
    \end{itemize}
    
    \textbf{D选项正确。}
    \begin{itemize}
        \item 经济资本计算：
        \[
        \begin{aligned}
        \text{经济资本} &= \text{敞口金额} \times \text{资本乘数} \times \sqrt{\text{PD标准差}^2 + \text{LR标准差}^2} \\
        &= 250 \times 7.50 \times \sqrt{0.08^2 + 0.15^2} \\
        &= 250 \times 7.50 \times 0.17 \\
        &= 45.00
        \end{aligned}
        \]
        \item 其中：
        \begin{enumerate}
            \item 敞口金额 = GBP 250 million
            \item 资本乘数 = 7.50
            \item PD标准差 = 8\%
            \item LR标准差 = 15\%
        \end{enumerate}
    \end{itemize}
\end{explanation}

\checklist{经济资本等于敞口金额乘以资本乘数再乘以风险标准差}


\begin{question}
    A team of credit analysts is using the Merton model to estimate the credit rating of a non-dividend-paying firm. The team compiles relevant information on the firm shown below:
    ·Value of total assets: AUD 250 million
    ·Long-term debt: AUD 120 million
    ·Short-term debt: AUD 0
    ·Expected return per year on firm assets: 12\%
    ·Asset volatility per year: 25\%
    The analysts use the following S\&P rating schedule provided by the investment bank to draw their conclusion

    \begin{table}[H]
        \centering
        \renewcommand{\arraystretch}{1.2}
        \begin{tabular}{|c|c|}
            \hline
            \textbf{Range of distance to default} & \textbf{S\&P rating class} \\
            \hline
            10.0 and greater & AAA \\
            \hline
            8.0 - 9.9 & AA \\
            \hline
            6.0 - 7.9 & A \\
            \hline
            4.0 - 5.9 & BBB \\
            \hline
            2.0 - 3.9 & BB \\
            \hline
            1.0 - 1.9 & B \\
            \hline
            0.5 - 0.9 & C \\
            \hline
        \end{tabular}
        \caption{距离违约区间与S\&P评级对应表}
    \end{table}

    What is the suggested credit rating of the firm at a 1-year horizon using the S\&P rating schedule?
\end{question}

\begin{choices}
    \item A
    \item BB
    \item B
    \item C
\end{choices}

\begin{correctanswer}
    B
\end{correctanswer}

\begin{explanation}
    \textbf{A选项错误。}
    \begin{itemize}
        \item 计算错误：
        \begin{enumerate}
            \item 可能低估了违约距离
            \item 可能使用了错误的参数
            \item 没有考虑资产波动率的影响
        \end{enumerate}
    \end{itemize}
    
    \textbf{B选项正确。}
    \begin{itemize}
        \item 违约距离计算：
        \[
        \begin{aligned}
        \text{违约距离} &= \frac{\ln\left(\frac{250}{120}\right) + (0.12 - 0.5 \times 0.25^2) \times 1}{0.25 \times \sqrt{1}} \\
        &= \frac{0.733 + 0.089}{0.25} \\
        &= 3.29
        \end{aligned}
        \]
        \item 根据S\&P评级表：
        \begin{enumerate}
            \item 违约距离 = 3.29
            \item 落在2.0 - 3.9区间
            \item 对应评级为BB
        \end{enumerate}
    \end{itemize}
    
    \textbf{C选项错误。}
    \begin{itemize}
        \item 计算错误：
        \begin{enumerate}
            \item 可能高估了违约距离
            \item 可能使用了错误的资产价值
            \item 可能忽略了预期收益率
        \end{enumerate}
    \end{itemize}
    
    \textbf{D选项错误。}
    \begin{itemize}
        \item 计算错误：
        \begin{enumerate}
            \item 可能使用了错误的波动率
            \item 可能低估了资产价值
            \item 可能高估了债务价值
        \end{enumerate}
    \end{itemize}
\end{explanation}

\checklist{使用Merton模型计算违约距离，根据S\&P评级表确定信用评级}



\newpage






\section{考点4:信用评分和零售信用风险管理}
\setcounter{question}{0}
\begin{question}
    Smith and Johnson divide credit analysis into four areas according to borrower type:
    Ⅰ.Consumer credit analysis is the evaluation of the creditworthiness of individual consumers;
    Ⅱ.Corporate credit analysis is the evaluation of nonfinancial companies such as manufacturers,
    and nonfinancial service providers; Ⅲ.Financial institution credit analysis is the evaluation
    of financial companies including banks and nonbank financial institutions, such as insurance
    companies and investment funds; Ⅳ.Sovereign/municipal credit analysis is the evaluation of the
    credit risk associated with the financial obligations of nations, subnational governments, and
    public authorities, as well as the impact of such risks on obligations of nonstate entities
    operating in specific jurisdictions. According to Smith and Johnson, each of the following is
    true about key features of credit analysis with respect to borrower type, except which is not
    true?
\end{question}

\begin{choices}
    \item Individuals (consumers): Credit analysis is amenable to automation and the use of scoring
    models and statistical tools to correlate risk to limited number of variables
    \item Non-financial corporations: Compared to consumers, tends to be more detailed and "hands-on"
    (i.e. less automated); key variables are likely to include liquidity, cash flow, near-term
    earnings capacity and profitability, solvency or capital position
    \item Financial Companies: In contrast to corporate (non-financial) credit analysis, qualitative
    analysis and asset quality are not important, but cash flow is a highly important (a "key
    indicator")
    \item Sovereigns: Includes analysis of country risk, which is primarily political dynamics and
    state of the economy; and systematic risk, which includes the regulatory regime and the
    financial system
\end{choices}

\begin{correctanswer}
    C
\end{correctanswer}

\begin{explanation}
    \textbf{A选项正确。}
    消费者信用分析适合自动化和使用评分模型，可以用统计工具将风险与有限数量的变量相关联。这是正确的，因为消费者信用分析通常涉及大量标准化数据，适合使用自动化工具和统计模型。

    \textbf{B选项正确。}
    非金融企业信用分析比消费者分析更详细且需要更多人工参与，关键变量包括流动性、现金流、盈利能力和资本状况。这是正确的，因为企业分析需要更全面的评估和更多定性因素。

    \textbf{C选项错误。}
    金融机构信用分析中，资产质量非常重要，而现金流不是关键指标。这与选项描述相反，因为金融机构的资产质量是评估其信用风险的关键因素，而现金流分析在金融机构信用评估中相对不那么重要。

    \textbf{D选项正确。}
    主权国家信用分析包括国家风险分析（政治动态和经济状况）和系统性风险分析（监管制度和金融体系）。这是正确的，因为主权信用分析需要同时考虑政治、经济和金融体系等多个维度。
\end{explanation}

\checklist{金融机构信用分析中资产质量重要而现金流不是关键指标}


\begin{question}
    Which of the following is not an example or element of predatory lending:
\end{question}

\begin{choices}
    \item Lender makes unaffordable loans based on borrower assets rather than ability to repay
    \item Lender induces borrower to repeatedly refinance in order to collect fees and charge high points
    \item Borrower misrepresents income or employment in mortgage application
    \item Lender engages in deception to conceal true nature of loan; e.g., deceives borrower into thinking loan is fixed-rate when mortgage is actually an adjustable-rate
\end{choices}

\begin{correctanswer}
    C
\end{correctanswer}

\begin{explanation}
    \textbf{A选项错误。}
    这是掠夺性贷款的一个典型特征。贷款人基于借款人的资产而不是还款能力发放贷款，这可能导致借款人无法负担贷款，属于掠夺性贷款行为。

    \textbf{B选项错误。}
    这是掠夺性贷款的另一个特征。贷款人诱导借款人反复再融资以收取费用和高额点数，这种行为被称为"贷款翻新"（loan flipping），是掠夺性贷款的常见手段。

    \textbf{C选项正确。}
    这是掠夺性借款（predatory borrowing）的例子，而不是掠夺性贷款（predatory lending）。借款人通过伪造收入或就业信息来获得贷款，这是借款人的欺诈行为，而不是贷款人的掠夺性行为。

    \textbf{D选项错误。}
    这是掠夺性贷款的第三个特征。贷款人通过欺骗手段隐瞒贷款的真实性质，例如将可调整利率贷款伪装成固定利率贷款，这属于掠夺性贷款行为。
\end{explanation}

\checklist{掠夺性借款是借款人欺诈而非贷款人掠夺性行为}


\begin{question}
    Which of the following was not a cause of the misalignment between investors, and rating agencies' incentives prior to the credit crisis of 2007-2009?
\end{question}

\begin{choices}
    \item Profit motive of the rating agencies
    \item Pressure from the originators of securitized products
    \item Manipulation of the ratings process by the originators
    \item Investors' lack of understanding of the products they were purchasing
\end{choices}

\begin{correctanswer}
    D
\end{correctanswer}

\begin{explanation}
    \textbf{A选项错误。}
    评级机构的利润动机是导致激励错位的原因之一。评级机构为了获取更多业务和利润，可能会降低评级标准，这导致了评级质量的下降。

    \textbf{B选项错误。}
    证券化产品发起人的压力是导致激励错位的重要原因。发起人通过向评级机构施压，要求获得更高的评级，这影响了评级的客观性。

    \textbf{C选项错误。}
    发起人对评级过程的操纵是导致激励错位的关键因素。发起人利用对评级机构评级流程的了解，设计出能够获得AAA评级的产品，这导致了评级的失真。

    \textbf{D选项正确。}
    投资者对产品缺乏了解不是导致激励错位的原因。根据金融危机调查委员会的调查结果，激励错位的主要原因是评级机构的模型缺陷、利润动机、发起人压力、市场份额驱动以及缺乏有效的公共监督等。
\end{explanation}

\checklist{投资者对产品缺乏了解不是评级机构激励错位的原因}

\begin{question}
    Which of the following statements is correct regarding credit risk scoring models?
\end{question}

\begin{choices}
    \item A pooled model will result in scores ranging from 300 to 850
    \item A custom model is cheaper to implement than credit bureau scores
    \item Multiple requests for new credit will reduce an applicant's credit score
    \item An example of a characteristic in a scoring model is the applicant's current gross salary of \$50,000
\end{choices}

\begin{correctanswer}
    C
\end{correctanswer}

\begin{explanation}
    \textbf{A选项错误。}
    信用局评分模型（而不是数据池模型）的分数范围是300到850。数据池模型使用相似的外部数据或内部数据，其分数范围可能不同。

    \textbf{B选项错误。}
    客户模型比信用局评分模型实施成本更高。客户模型需要银行自己建立和维护，而信用局评分模型可以直接使用外部专业机构提供的服务。

    \textbf{C选项正确。}
    多次申请新信用会降低申请人的信用评分。信用记录会显示信用申请历史，频繁的申请会被视为负面因素，导致信用评分下降。

    \textbf{D选项错误。}
    "当前雇主的总工资"是一个特征（characteristic），而具体的工资数额（\$50,000）是一个属性（attribute）。在评分模型中，特征是指变量类别，而属性是指具体的数值。
\end{explanation}

\checklist{多次申请新信用会降低申请人的信用评分}


\begin{question}
    At the beginning of the year, a firm bought an AA-rated corporate bond at USD 115 per USD 100 face value. Using market data, the risk manager estimates the following year-end values for the bond based on interest rate simulations informed by the economics team:

    \begin{table}[htbp]
        \centering
        \begin{tabular}{|c|c|}
            \hline
            Rating & Year-end Bond Value \\
            & (USD per USD 100 face value) \\
            \hline
            AAA     & 118 \\
            AA      & 114 \\
            A       & 108 \\
            BBB     & 103 \\
            BB      & 95  \\
            B       & 85  \\
            CCC     & 75  \\
            Default & 45  \\
            \hline
        \end{tabular}
        \caption{不同评级债券年末价值}
    \end{table}

    In addition, the risk manager estimates the 1-year transition probabilities on the AA-rated corporate bond: 
    \begin{table}[htbp]
        \centering
        \begin{tabular}{|c|c|}
            \hline
            Rating & Probability of State \\
            \hline
            AAA     & 4.00\% \\
            AA      & 82.00\% \\
            A       & 8.00\% \\
            BBB     & 4.50\% \\
            BB      & 0.40\% \\
            B       & 0.30\% \\
            CCC     & 0.20\% \\
            Default & 0.60\% \\
            \hline
        \end{tabular}
        \caption{各评级状态概率}
    \end{table}

    What is the 1-year 95\% credit VaR per USD 100 of face value closest to?
\end{question}

\begin{choices}
    \item USD 12
    \item USD 20
    \item USD 35
    \item USD 45
\end{choices}

\begin{correctanswer}
    A
\end{correctanswer}

\begin{explanation}
    \textbf{A选项正确。}
    95\%信用VaR等于当前价值减去95\%分位数下的预期价值。根据评级转移概率表，95\%分位数对应BBB评级，此时债券价值为103。因此，信用VaR = 115 - 103 = 12。

    \textbf{B选项错误。}
    这个答案可能错误地使用了BB评级对应的价值（95）来计算VaR，导致结果过大。

    \textbf{C选项错误。}
    这个答案可能错误地使用了B评级对应的价值（85）来计算VaR，导致结果过大。

    \textbf{D选项错误。}
    这个答案可能错误地使用了违约状态对应的价值（45）来计算VaR，导致结果过大。
\end{explanation}

\checklist{信用VaR等于当前价值减去95\%分位数下的预期价值}


\begin{question}
    A risk analyst at a bank is preparing a presentation on climate-related financial risks. The analyst focuses on the amplifiers and mitigants of climate-related financial risks and assesses the manner in which they function. Which of the following statements would be correct for the analyst to include in the presentation?
\end{question}

\begin{choices}
    \item Climate-related risks increase through asset securitizations
    \item Climate-related risks are typically impossible to hedge
    \item Physical and transition risk drivers interact with each other
    \item Transition risk drivers generally affect firms idiosyncratically
\end{choices}

\begin{correctanswer}
    C
\end{correctanswer}

\begin{explanation}
    \textbf{A选项错误。}
    资产证券化并不会增加气候相关风险。资产证券化是一种风险分散和转移的工具，它本身不会放大气候风险。

    \textbf{B选项错误。}
    气候相关风险并非完全无法对冲。虽然完全对冲可能很困难，但通过保险、衍生品等工具可以部分对冲这些风险。

    \textbf{C选项正确。}
    物理风险和转型风险驱动因素会相互影响。例如，气候变化导致的物理风险可能加速向低碳经济的转型，而转型政策又可能影响物理风险的程度。

    \textbf{D选项错误。}
    转型风险驱动因素通常会影响整个行业或经济系统，而不是仅仅影响个别公司。例如，碳定价政策会影响所有高碳排放企业。
\end{explanation}

\checklist{物理风险和转型风险驱动因素会相互影响}





\newpage


\section{考点5:主权国家违约风险}
\setcounter{question}{0}


\begin{question}
    You are currently long \$12,000,000 par value, 7% XYZ bonds. To hedge your position, you must decide between credit protection via a 5-year CDS with 70bp annual premiums or digital swap with 55% payout with 60bp annual premiums. After one year, XYZ has defaulted on its debt obligations and currently trades at 65% of par. Which of the following statements is true?
\end{question}

\begin{choices}
    \item The contingent payment from the protection buyer to the protection seller is greater under the single-name CDS than the digital swap
    \item The contingent payment from the protection buyer to the protection seller is less under the single-name CDS than the digital swap
    \item The contingent payment from the protection seller to the protection buyer is greater under the single-name CDS than the digital swap
    \item The contingent payment from the protection seller to the protection buyer is less under the single-name CDS than the digital swap
\end{choices}

\begin{correctanswer}
    D
\end{correctanswer}

\begin{explanation}
    \textbf{A选项错误。}
    违约时，赔付是从保护卖方支付给保护买方，而不是从买方支付给卖方。因此这个选项描述的方向是错误的。

    \textbf{B选项错误。}
    同样，这个选项也错误地描述了赔付方向。违约时，赔付是从保护卖方支付给保护买方。

    \textbf{C选项错误。}
    实物交割CDS的赔付金额是面值的35\%（1-0.65），而数字互换的赔付金额是面值的55\%。因此，数字互换的赔付金额更大。

    \textbf{D选项正确。}
    实物交割CDS的赔付金额是面值的35\%（1-0.65），而数字互换的赔付金额是面值的55\%。因此，实物交割CDS的赔付金额小于数字互换。
\end{explanation}

\checklist{实物交割CDS的赔付金额小于数字互换的赔付金额}


\begin{question}
    A risk manager is advising the trading desk about entering into a digital credit default swap as a way to obtain credit protection. Which cash flow and delivery requirement will the desk most likely experience in the event of a default of the underlying reference asset?
\end{question}

\begin{choices}
    \item Receive the pre-agreed cash payment; delivering nothing
    \item Receive [(Par Value) – (Market Value of Reference Asset)]; deliver the reference asset
    \item Receive [(Par Value) – (Market Value of Reference Asset)]; deliver nothing
    \item Receive the pre-agreed cash payment; deliver the reference asset
\end{choices}

\begin{correctanswer}
    A
\end{correctanswer}

\begin{explanation}
    \textbf{A选项正确。}
    数字式信用违约互换（Digital CDS）在违约时会支付预先约定的固定金额，且不需要交付任何资产。这种设计特别适用于流动性差、难以定价的参考资产。

    \textbf{B选项错误。}
    这个选项描述的是实物交割CDS的特点，而不是数字式CDS。实物交割CDS需要交付参考资产，并支付面值与市场价值的差额。

    \textbf{C选项错误。}
    这个选项描述的是现金结算CDS的特点，而不是数字式CDS。现金结算CDS支付面值与市场价值的差额，但不需要交付资产。

    \textbf{D选项错误。}
    数字式CDS不需要交付参考资产，只需要支付预先约定的固定金额。这个选项错误地要求交付参考资产。
\end{explanation}

\checklist{数字式CDS在违约时支付固定金额且无需交付资产}


\begin{question}
    A default swap acts like a:
\end{question}

\begin{choices}
    \item Call option on the reference obligation for the buyer of the swap
    \item Put option on the reference obligation for the buyer of the swap
    \item Look back option on the reference obligation for the buyer of the swap
    \item Up-and-out option on the reference obligation for the buyer of the swap
\end{choices}

\begin{correctanswer}
    B
\end{correctanswer}

\begin{explanation}
    \textbf{A选项错误。}
    看涨期权赋予买方以固定价格购买资产的权利，而违约互换是保护买方免受违约损失的工具，两者性质不同。

    \textbf{B选项正确。}
    违约互换的行为类似于看跌期权。当参考资产违约时，买方获得赔付，这限制了买方的下行风险，这与看跌期权的保护功能相似。

    \textbf{C选项错误。}
    回望期权允许买方在期权有效期内选择最优执行价格，而违约互换的赔付是预先确定的，两者机制不同。

    \textbf{D选项错误。}
    向上敲出期权在资产价格超过特定水平时自动失效，而违约互换在违约发生时才会触发赔付，两者触发条件不同。
\end{explanation}

\checklist{违约互换的行为类似于看跌期权}

\begin{question}
    In comparing agency rating systems to internal (experts-based) rating systems, evidence has shown that:
\end{question}

\begin{choices}
    \item Internal systems and agency systems are equally compliant in regard to specificity
    \item Agency systems are more compliant in regard to verifiability than internal systems
    \item Agency systems are less compliant in regard to measurability than internal systems
    \item Internal systems are more compliant in regard to objectivity and homogeneity than agency systems
\end{choices}

\begin{correctanswer}
    B
\end{correctanswer}

\begin{explanation}
    \textbf{A选项错误。}
    内部评级系统和机构评级系统在特异性方面并不等同。机构评级系统通常有更明确的评级标准和更一致的评级方法。

    \textbf{B选项正确。}
    机构评级系统在可验证性方面比内部评级系统更符合要求。机构评级系统有更透明的评级过程和更严格的验证机制，使得评级结果更容易被验证。

    \textbf{C选项错误。}
    机构评级系统在可测量性方面并不比内部评级系统差。相反，机构评级系统通常有更完善的量化指标和更系统的测量方法。

    \textbf{D选项错误。}
    内部评级系统在客观性和同质性方面并不比机构评级系统更好。机构评级系统通常有更严格的评级标准和更统一的评级方法，能够提供更客观和一致的评级结果。
\end{explanation}

\checklist{机构评级系统在可验证性方面优于内部评级系统}

\begin{question}
    The buyer of a credit-default swap (CDS):
\end{question}

\begin{choices}
    \item Has no risk exposure to the combined holdings of the CDS and reference obligation
    \item Is subject to the credit risk of the seller
    \item Must place collateral with the seller of the CDS to ensure performance
    \item Can surrender the CDS for the value of the accumulated payments
\end{choices}

\begin{correctanswer}
    B
\end{correctanswer}

\begin{explanation}
    \textbf{A选项错误。}
    CDS买方仍然面临风险敞口。即使持有CDS和参考资产，买方仍然面临卖方的信用风险，无法完全消除风险。

    \textbf{B选项正确。}
    CDS买方面临卖方的信用风险。如果参考资产违约，而卖方无法履行支付义务，买方将遭受损失。这种风险受两个因素影响：1）参考资产违约概率随时间增加，卖方违约概率也可能增加；2）经济和运营因素可能导致参考资产违约与卖方违约之间的相关性增加。

    \textbf{C选项错误。}
    CDS买方不需要向卖方提供抵押品。相反，通常是卖方需要提供抵押品来确保其履行支付义务的能力。

    \textbf{D选项错误。}
    CDS买方不能通过放弃CDS来获得累积付款的价值。CDS是一种信用保护工具，不能像债券那样在二级市场交易或赎回。
\end{explanation}

\checklist{CDS买方面临卖方的信用风险}

\begin{question}
    A quantitative risk analyst at a broker-dealer bank is studying correlations between financial assets, using historical data over the last twenty years, with particular focus on the 2007-2009 global financial crisis. The analyst also examines the tools and products used to gain exposure to or hedge against changes in correlations. Which of the following can the analyst correctly conclude?
\end{question}

\begin{choices}
    \item The spreads between prices of CDO tranches will decrease as mortgage default probabilities increase during a systemic housing market crash
    \item Both VaR and ES models incorporate the impacts of dynamic financial correlations between assets to manage risk and estimate capital requirements
    \item The correlation swap is the only financial product available to market participants to hedge or expand their exposure to the correlation between assets
    \item The CDS spread would decrease if the default correlation between the reference asset and the seller of the CDS on that asset increases
\end{choices}

\begin{correctanswer}
    D
\end{correctanswer}

\begin{explanation}
    \textbf{A选项错误。}
    在系统性房地产市场崩溃期间，随着抵押贷款违约概率增加，CDO分层的价格利差会扩大而不是缩小。这是因为违约风险增加会导致各层级的风险溢价上升。

    \textbf{B选项错误。}
    VaR和ES模型并不总是能够充分捕捉资产之间的动态金融相关性。这些模型通常假设相关性是稳定的，而实际上相关性在危机期间会发生显著变化。

    \textbf{C选项错误。}
    相关性互换不是唯一可用于对冲或扩大资产相关性敞口的金融产品。市场参与者还可以使用其他工具，如CDO、CDS指数等。

    \textbf{D选项正确。}
    如果参考资产与CDS卖方之间的违约相关性增加，CDS利差会下降。这是因为违约相关性增加意味着卖方更可能无法履行支付义务，从而降低了CDS的价值。
\end{explanation}

\checklist{参考资产与CDS卖方违约相关性增加会导致CDS利差下降}

\newpage



\end{document}